% ---------------------------------------------------------------------------
% transfer_ext.tex: transfers and extension-field expansion.
% Five chapters, to sit after ch:extension (sections/extension.tex) and
% sec:transfer (sections/prog_transfer.tex). Draft of 2026-11-06.
% Every figure is producer-reported (an unverified assertion in the sense of
% the proof-carrying-claims rule) unless the text names a passed validation
% gate (a VALIDATION.json with a passing validator report) for it.
% ---------------------------------------------------------------------------

\chapter{Isogenies, volcanoes and conductor gaps}\label{ch:isogenies}

Chapter~\ref{sec:transfer} states the transport-conservation lemma
(Lemma~\ref{lem:conservation}) and records the volcano census, the
large-sample floor/crater ratio, the ECC2K-130 conductor-gap rounds and
fibre-aware index calculus. Chapter~\ref{ch:obstructions} draws the
consequence for factor bases: an isogeny buys only the kernel symmetry,
the destination's factor-base geometry and its arithmetic. This chapter
collects what the programme measured or proved on the isogeny side beyond
those summaries: the closure that makes a prime-field isogeny walk to a
weak curve impossible, the index-calculus readings across volcano levels,
the conductor-gap constructions and their correction trail, and the
factor bases and seeds that use isogenies. The rho side of the same
studies (one instance on every level, whole classes, the ecbench volcano
sessions) is in the matched-reference chapters (\S\ref{sec:rv-isogeny});
it is not repeated here. Unless a validation gate is named, every number in
this chapter is a producer-reported figure without a pcc standing.

\section{The Tate closure on prime fields}\label{sec:iso-tate}

\paragraph{The statement.} Two ordinary elliptic curves over $\FF_p$ are
$\FF_p$-isogenous if and only if they have the same number of rational
points (Tate's isogeny theorem). The three weak families an isogeny walk
might hope to reach are all functions of $(\#E(\FF_p),p)$:
\begin{itemize}
\item an anomalous curve has $\#E(\FF_p)=p$, trace one;
\item the embedding degree is $k=\operatorname{ord}_N(p)$ for the prime
  $N\mid\#E$;
\item a GHS or Weil descent needs a proper subfield, and $\FF_p$ has none.
\end{itemize}
So the weak set inside the isogeny class of a generic prime-field curve is
empty at every size, not merely far away. The programme filed this as a
derivation, strength \emph{strong}, against its own hypothesis
\texttt{H-IT-001} (\art{crypto-autoresearcher:ledger/evidence/EV-IT-511f3d.yaml}).
The record notes that an earlier evidence item had stated the obstruction as
a path length of at least $p/\ell^2$; the theorem gives disconnection, which
is stronger. The decision \texttt{DEC-20260807-d4f2a9} rejected the lead in
scope; \texttt{DEC-20260808-86ec3c} later restored its status to
``weakened'' under the programme's no-pausing directive. That second
decision is administrative: it adds no evidence
(\art{crypto-autoresearcher:knowledge/leads/ecdlp/transfer/LEAD-TRF-prime-field-isogeny-to-weak.yaml}).

\paragraph{The measured null.} Before the closure was filed, experiment
\texttt{EXP-IT-001} searched for special curves reachable from generic ones
and found none: the special-curve share $\rho_{\rm special}$ was $0$ at
20 bits (exhaustive), at 24 bits (78\% coverage, stopped on resources) and
at 28 bits (a 0.16\% sample). Its positive control was a Smart certificate
on an anomalous curve (\S\ref{sec:tp-smart}). Both are in
\art{crypto-autoresearcher:ledger/evidence/EV-IT-aefd12.yaml}. The scope is
prime fields. Over composite-degree binary fields descent weakness is
\emph{not} an isogeny-class invariant, which is what makes Teske's trapdoor
and the walks of Chapter~\ref{ch:ghs} possible.

\paragraph{Walks of standard curves.} Several campaigns walked the
isogeny graphs of the NIST prime curves as tooling, without an ECDLP cost.
\begin{itemize}
\item From registered P-256, the native walker produced 65{,}536 distinct
  curve identities and 302{,}653 certified directed isogeny edges across the 12 walkable odd degrees up to 61; every order was proved prime
  and every model nonsingular with a valid generator. The artefacts were
  not published to S3, so the gate passed locally only
  (\art{research/p256_isogeny_multigraph_64k_20261006/RESULTS.md}).
\item On P-192 and P-224, 100 distinct large prime-degree maps were built
  and replayed independently: 50 consecutive edges of degree 2341 in the
  P-192 class and 50 of degree 1117 in the P-224 class, with exact kernel,
  codomain, equation, subgroup-order and scalar-transport checks
  (\art{research/p192_p224_hundred_large_prime_isogenies_20261010/RESULTS.md}).
  The screen behind it covered every prime from 1010 to $10^6$ on both
  sources.
\end{itemize}
None of these maps can change the ECDLP cost: by Tate and
Lemma~\ref{lem:conservation} the destination has the same group and the
same factor-base options up to arithmetic.

\paragraph{A battery on the NIST $b$ coefficients.} A frozen protocol
tested whether the published $b$ of five standard curves is unusual under
17 statistics (\art{research/b_weak_class_battery_20261022/README.md}).
The rule was Bonferroni over $m=5\times17=85$ tests, flagging only
$p\le5.88\cdot10^{-4}$. No curve--statistic pair was flagged; the smallest
$p$ was $0.034$ (P-256, statistic \texttt{smooth\_c}), 57 times the
threshold. Positive controls passed in 74 of 85 cells and 11 cells were
inconclusive. This is a negative-control result about representation, not
about the isogeny class.

\paragraph{Designs that did not run to an answer.} The class-wide yield
study for the ECCp-131 and ECC2K-130 classes
(\art{research/isogeny_class_ic_yield_20261102/}) has a design and run
cards; one ECC2K-130 card timed out at 1{,}872\,s with 1{,}713 of 2{,}112
rows written, because an untimed solution count grew by a factor of 8 per
step of $k$ (\texttt{TIMEOUT\_NOTE.md}). Its partial output is not analysed
and is not cited here.

\section{Index calculus across volcano levels}\label{sec:iso-volcano}

Within one volcano every curve has the same $\#E$, so generic cost is
level-blind by the generic-group argument, and the measured rho nulls of
\S\ref{sec:rv-isogeny} agree. Index calculus is not generic: its factor
base is a set of points chosen by a coordinate rule, so it could in
principle see the level, the endomorphism ring or the model. Four studies
measured that, and one directory must be excluded.

\paragraph{Vertical, horizontal and coordinate effects.} The frozen
volcano-cost study ran rho, BSGS and index calculus on six 24-bit classes
of depth 3 ($\ell\in\{2,3\}$), 200 curves in four isomorphic models each:
43{,}200 measured runs, 46{,}912 with the two control classes, every run
verified and 408 of 408 replays exact
(\art{research/isogeny-volcano-cost_20261017/README.md}).
Table~\ref{tab:iso-volcano-v1} gives the index-calculus rows under the
v1.1 statistic. BSGS cost was identical on every curve and model.

\begin{table}[t]
\centering\small
\begin{tabularx}{\linewidth}{@{}>{\raggedright}p{2.6cm}>{\raggedright}X>{\raggedright}X>{\raggedright\arraybackslash}X@{}}
\toprule
Metric & Vertical (floor/crater, 90\%) & Horizontal $\tau$ (null 95th pct.), $p$ & Coordinates $\tau$ (null 95th pct.), $p$ \\
\midrule
rho \texttt{rho.negation} $S$ & $1.022$ $[0.983,1.062]$ & $1.7\%$ ($2.6\%$), $0.28$ & $2.6\%$ ($4.6\%$), $0.70$ \\
IC relation yield & $0.980$ $[0.957,1.004]$; inconclusive & $2.2\%$ ($1.4\%$), $0.0005$ & $4.7\%$ ($2.5\%$), $0.0005$ \\
IC cost $S$ (reported) & $1.039$ $[1.008,1.071]$ & $2.9\%$ ($2.1\%$), $0.002$ & $5.6\%$ ($3.7\%$), $0.0005$ \\
\bottomrule
\end{tabularx}
\caption[]{Volcano-level effects on index calculus (\texttt{ic.pipeline},
the 32 smallest abscissae, meet-in-the-middle) at 24 bits, protocol v1.1.
$\tau$: dispersion between curves of one level (horizontal) or between
isomorphic models of one curve (coordinates). The vertical IC-yield test
is inconclusive because its planted $+5\%$ floor control was not detected.
Source: \art{research/isogeny-volcano-cost_20261017/README.md}.}
\label{tab:iso-volcano-v1}
\end{table}

Three readings follow, each stated by the source.
\begin{itemize}
\item \emph{The coordinate effect is larger than the isogeny effect.}
  Isomorphic models of one curve differ in relation yield by
  $\tau\approx4.7\%$, more than curves of one level do ($2.2\%$). The
  factor base ``the 32 points of smallest abscissa'' is a different point
  set in each model, and the yield follows the set, not the curve. By the
  protocol's own rule this is a coordinate effect.
\item \emph{The v1 statistic was wrong, and the report says so.} Under the
  frozen v1 statistic the vertical IC-yield test was nominally significant
  ($0.972$ $[0.951,0.993]$, $p=0.0012$). The same v1 statistic gave BSGS,
  whose cost is identical on every curve, a bootstrap interval of
  $[0.88,1.22]$, because it pools classes with different shares of floor
  and crater curves. The amended v1.1 statistic forms the ratio inside each
  class; under it the test is not significant. The study does not claim the
  v1 reading of a $-2.8\%$ level effect.
\item \emph{The distance to rho is untouched.} Per curve, index calculus
  costs $S\approx27$--$96$ at $w=0$ against rho's $0.8$--$1.6$, so it is
  about 30--60 times rho at 24 bits, and rho is below index calculus on
  every one of the 800 workload curves at every weight of the band.
\end{itemize}
An exploratory per-class contrast (six unadjusted contrasts) suggests that
on $\ell=2$ volcanoes floor curves, whose 2-part is cyclic, give about
$6.5\%$ fewer relations per trial than crater curves with full rational
2-torsion ($0.930$ $[0.863,0.994]$ and $0.937$ $[0.888,0.987]$ on the
ramified and split classes), while $\ell=3$ volcanoes show nothing. The
source files this as a hypothesis for a later protocol. Its successor
attributed $99\%$ of the $\ell=2$ contrast to the base's cofactor classes
(Chapter~\ref{sec:transfer}). The study's crater controls needed the
24 October class-set-up erratum, which is in Table~\ref{tab:corr-rho}.

\paragraph{Tall volcanoes.} A preregistered study on ten volcanoes with
$\ell\in\{2,3,5,7\}$ and depths 6 to 24, whose tallest floors hold 8.4 and
9.6 million curves, ran 14{,}731 \texttt{rho.negation} solves on 608
checked curves (\art{research/volcano_tall_layers_20261018/RESULT.md}).
Every preregistered comparison lies inside the $\pm5\%$ margin: in tier a
(eight volcanoes) floor/crater is $1.0070$ $[0.9741,1.0394]$, middle/crater
$0.9898$ $[0.9654,1.0139]$, and the crater-to-floor slope $+0.87\%$
$[-2.03\%,+3.75\%]$; the Kolmogorov--Smirnov test between crater and floor
gives $D=0.0199$, $p=0.715$. The positive control, plain over negation rho
on the same workloads, reads $1.3485$ $[1.2380,1.4780]$, so the design
resolves a 35\% difference. None of the 40 per-volcano intervals excludes
the null; the widest are $1.079$ $[0.973,1.188]$ (a8, $\ell=7$, depth 6) and
$0.961$ $[0.879,1.050]$ (a7). The raw $S$ of negation rho is $0.947$
$[0.936,0.958]$ against the floor $\sqrt{\pi/4}=0.8862$, constant from
$r\approx2^{42}$ to $2^{51}$. Three tier-b sessions were cut short by a
worker restart and kept as measured. The study retracts an earlier
``$\pm8\%$'' quoted for the Koblitz layers below.

\paragraph{Prime and Koblitz layers.} Two earlier studies have wide
intervals and say so.
\begin{itemize}
\item On two depth-8 $\ell=2$ volcanoes over $\FF_p$ (864 records,
  \art{research/volcano_prime_layers_20261018/RESULT.md}), 14 of 16 tested
  intervals contain the null. The two that exclude it are both
  \texttt{rho.plain} on one volcano ($m_5$: middle/crater $0.688$
  $[0.562,0.891]$, floor/crater $0.671$ $[0.524,0.883]$), and both involve
  its crater, a single curve with six targets, so a crater effect cannot be
  separated from a curve effect. The non-excluding intervals are about
  $\pm30\%$ wide.
\item On binary Koblitz curves with $n\in\{19,23,37,41,61\}$, level-1 over
  crater for \texttt{rho.negation} contains 1 at every $n$, from $0.837$
  $[0.625,1.073]$ ($n=61$) to $1.089$ $[0.993,1.195]$ ($n=19$); one of ten
  plain-rho intervals excludes 1 ($n=41$, $1.152$ $[1.021,1.304]$). The
  Frobenius solver over the negation solver on the crater falls from
  $2.153$ ($n=19$, fixed overhead) to $0.150$ $[0.132,0.170]$ at $n=41$,
  against the ideal $1/\sqrt n=0.156$
  (\art{research/volcano_koblitz_layers_20261018/RESULT.md}). The
  $\sqrt n$ gain is the crater's automorphism, which its neighbours lose.
\item At $n=83$ the native tool built the modular polynomial
  $\Phi_{6473}(X,j_0)$ from 9{,}864 CRT primes in 8{,}590\,s, reduced and
  factored it into 78 Frobenius orbits of size 83, and found the expected
  6{,}474 level-1 $j$-invariants; no ECDLP was measured there
  (\art{research/volcano_koblitz_n83_20261018/}).
\end{itemize}

\paragraph{An excluded directory.} \art{research/synthetic_volcano_ecdlp_20261018/}
states in its first line that its data are invented: no elliptic curve,
point, isogeny map, relation matrix or logarithm was constructed, and its
130{,}410 ``observations'' come from a planning model. It is not evidence
for any statement in this treatise and is not cited for one.

\section{Conductor gaps}\label{sec:iso-conductor}

Jao, Miller and Venkatesan~\cite{jmv2005} leave open whether curves of one
order whose endomorphism rings differ by a large conductor gap could differ
in hardness. Chapter~\ref{sec:transfer} gives the ECC2K-130 rounds and the
fibre-scaling law; the corrections chapter gives the sheriff findings on the
CM-pairs banner and the round-2 set-up defect (Table~\ref{tab:corr-structure},
\S\ref{sec:corr-structure}). This section records the constructions and the
index-calculus measurements those summaries omit.

\paragraph{Certifying a gap without walking it.} With a class-number-one CM
field the gap can be certified from Frobenius alone
(\art{research/conductor_gap_cm_pairs_20261010/README.md} \S1). Take
$D=-19$, a prime $\ell$, and odd $t$ with $q=(t^2+19\ell^2)/4$ and
$N=q+1-t$ both prime. Then $t^2-4q=-19\ell^2$, so every ordinary curve of
trace $t$ has conductor $1$ or $\ell$; the curve with $j=-884736$ is the
crater and every other curve of order $N$ is a floor curve. Floor curves are
found by sampling $(a,b)$ and testing $[N]P=O$, about $19\ell/4$ samples per
hit; no isogeny is computed, and for $\ell\ge263$ nothing in the repository
could exhibit one. The construction fails for $D=-7$ and $D\equiv1\pmod8$
(conductor forced even) and for $D=-11$ ($3\mid qN$ for every odd $t$);
$D\in\{-19,-43,-67,-163\}$ work. Pairs were built at $\ell=263$ (19-bit
$q$), $8209$ (29 bits) and $131101$ (37 bits) and checked independently in
Sage (Table~\ref{tab:iso-gap-pairs}).

\begin{table}[t]
\centering\small
\begin{tabular}{@{}rrlrr@{}}
\toprule
$\ell$ & $q$ bits & curves & samples per floor hit & predicted \\
\midrule
263 & 19 & 1 crater, 4 floor & $\approx1.3$k & $\approx1.2$k \\
8209 & 29 & 1 crater, 3 floor & $\approx96$k & $\approx39$k \\
131101 & 37 & 1 crater, 1 floor & $\approx970$k & $\approx620$k \\
\bottomrule
\end{tabular}
\caption[]{Certified conductor-gap CM pairs with $D=-19$: crater $j=-884736$,
floor curves found by sampling and one scalar multiplication per sample.
Observed rates agree with the prediction within the Poisson noise of one to
three hits. Every curve was re-checked in Sage: $\#E=N$ by SEA, the
$j$-invariant, the Frobenius discriminant $-19\ell^2$ and an embedding degree
above 20. Source: \art{research/conductor_gap_cm_pairs_20261010/README.md}
\S\S2--3.}
\label{tab:iso-gap-pairs}
\end{table} The study's banner was the one sheriff cases
\texttt{SHF-20261017-e4d7f6} and \texttt{SHF-20261028-0f3f64} withdrew and
then corrected: it now reads that no implemented method shows a
\emph{reproducible} conductor effect, pending a rho re-test. Its Lean
generic-invariance theorem covers only algorithms whose control flow depends
on equality answers and a random tape, that is BSGS-type algorithms; the
rho and kangaroo that ecbench runs hash point representations and are
outside it (\S4 of the source, corrected 16 October).

\paragraph{Fibre-aware relation generation across the gap.} The
index-calculus study of the same date built 38 prime-field pairs with gaps
121 to 1458, five Koblitz pairs and three random binary pairs, and three
fibre-aware oracles: \emph{lifts} (probe every cofactor lift of the
target), \emph{closed} (a pair table holding $P_i+P_j+K$ for every
cofactor-torsion $K$) and \emph{combined} (one Weil-descended system with
the fibre polynomial adjoined)
(\art{research/isogeny_conductor_gap_ic_20261010/README.md}).
Table~\ref{tab:iso-fibre} gives Part~A.

\begin{table}[t]
\centering\small
\begin{tabular}{@{}llrrr@{}}
\toprule
Family, cofactor & Arm & $S$/blind & yield/blind & cost per relation/blind \\
\midrule
prime, $h\le25$ & closed & $0.50$ $[0.21,2.36]$ & $11.7$ & $0.094$ \\
 & lifts & $1.83$ $[1.00,3.02]$ & $11.7$ & $1.01$ \\
prime, $h\ge50$ & closed & $4.7$ $[2.3,75]$ & $116$ & $0.017$ \\
 & lifts & $3.8$ $[1.8,29]$ & $114$ & $0.86$ \\
\addlinespace
random binary, $h\in\{4,6,8\}$ & closed MITM & $5.8$ $[3.9,7.9]$ & $3.5$ & $0.65$ \\
 & lifts MITM & $1.02$ & $3.5$ & $0.99$ \\
Koblitz, $h\in\{2,4\}$ & closed MITM & $3.0$ $[1.9,4.0]$ & $2.2$ & $0.78$ \\
 & lifts MITM & $1.01$ & $2.2$ & $1.04$ \\
\bottomrule
\end{tabular}
\caption[]{Fibre-aware arms against the blind oracle on the same curve,
factor base and targets (medians, with the range for $S$). $S$ is the cold
whole-pipeline cost over $\sqrt r$, a lower bound. Source:
\art{research/isogeny_conductor_gap_ic_20261010/README.md} \S3.1.}
\label{tab:iso-fibre}
\end{table}

Yield scales with the fibre, up to the cofactor $h$, as predicted. Only the
closed table cuts the relation phase's cost per relation, and it pays in
set-up: it wins on prime curves with $h\le25$ ($S$ falls to $0.21$--$0.36$
of blind on the 25-bit curves) and loses everywhere else, up to $75\times$
blind at $h=784$. The F4 solving degree does not move for lifts (mean
$2.96$--$3.00$). No fibre arm beats rho: the best cell, closed MITM on the
25-bit curves, has $S\ge21$ against rho's $1.0$--$1.4$. The combined Weil
descent is negative, measured on warm-up runs only, because host reboots
ended all four attempts at its session. The study is classified as
engineering on the relation phase of cofactor curves; the conductor plays
no role in any of these ratios. The next round's correction A0 (a
floor/crater set-up ratio of $3.10$ that was a counting defect, corrected to
$1.00$) is in Table~\ref{tab:corr-structure}.

\paragraph{Drastic gaps.} Two rounds pushed the gap to $\ell\approx2^{19}$
to $2^{21}$, 4 to 16 times the largest earlier gap.
\begin{itemize}
\item Round~1 built six certified pairs with $D=-19$ and $D=-43$ at
  $\ell\in\{508433,1016839,2033609\}$, eight one-target workloads per
  curve (\art{research/drastic_prime_conductor_gap_20261012/README.md}).
  The rho floor/crater ratios range from $0.529$ to $1.293$; only one 95\%
  interval excludes 1 ($D=-43$, $\ell=1016839$: $0.529$
  $[0.306,0.874]$, permutation $p=0.038$), and its 99\% interval does not.
  Index calculus was not run at these sizes.
\item Round~2 took the $j=0$ crater of $D=-3$, whose automorphism group has
  order 6, at $\ell=1016839$ and $2033609$ with 32 workloads per curve
  (\art{research/drastic_prime_conductor_gap_round2_20261012/README.md}).
  The same-method rho ratio sits at $0.947$ and $1.060$; the best-method
  ratio (\texttt{rho.cubic} on the crater, negation rho on the floor) is
  $1.692$ $[1.217,2.362]$ and $1.730$ $[1.247,2.461]$ at 99\%, within 3\%
  of $\sqrt3=1.7321$. The round-1 lead at $D=-43$ re-measured on 32 fresh
  targets at $0.947$ $[0.662,1.347]$: target-sampling noise.
\end{itemize}
The $\sqrt3$ is the automorphism group, not the conductor. Round~3's
confirmation and its later set-up correction are in
Chapter~\ref{sec:transfer} and \S\ref{sec:rv-aut}.

\paragraph{CM orbits and a generic recurrence.} The 15 October study ran an
orbit-folded index calculus and rho on 100 gap pairs
(\art{research/conductor_gap_cm_orbits_20261015/REPORT.md}). All 500
primary attempts verified (300 IC, 200 rho). Floor/crater native CPU for
index calculus is $0.980$ $[0.866,1.111]$ and the query-count ratio
$0.878$ $[0.764,1.009]$; neither establishes a difference. A post-hoc
paired rho diagnostic gives CPU $1.174$ $[0.999,1.383]$ and charged work
$1.160$ $[0.986,1.367]$. The replication cohort shows how fragile such
readings are: fresh targets give $0.988$ $[0.845,1.149]$ in queries, but
isomorphic coordinates give $1.397$ $[0.882,2.056]$ and $D=-43$ gives
$1.375$ $[0.998,2.011]$, whose CPU interval narrowly excludes 1 in an
unadjusted exploratory subset. All timings are co-scheduled L0, with crater
and floor shards in mostly separate queue waves, so the source calls them
descriptive. The study also found that its baseline was too weak: the CM
action satisfies $\lambda^2-\lambda+5=0$ modulo the group order, so the
generic recurrence $P_{k+2}=P_{k+1}-[5]P_k$ generates the same orbit on
either endpoint without a degree-5 endomorphism; per 8{,}192-point
construction the crater CM map costs $5.265$\,ms, either recurrence
$6.64$\,ms and a direct floor scalar multiplication $41.494$\,ms. The
comparison against repeated full scalar multiplication overstated the
advantage the crater's endomorphism offers. A 96-bit rational-family floor
was certified separately; ECDLP timings remain at 53 bits.

\paragraph{Follow-ups and archives.} The five follow-up experiments of
16 October (interleaved crossover, relation-yield distribution, orbit
geometry, conductor-size trend, rho partition controls) were frozen and
launched; at the published checkpoint 182 crossover receipts existed (178
recoveries, four wall-censored), which the source calls execution status,
not a comparison (\art{research/conductor_gap_followups_20261016/README.md}).
Rounds 5--7 of the ECC2K-130 series, the replication, size-ladder and
crater-design directories, and the 60- and 80-bit archival imports have
records or PDFs but no README result statement and no admitted row; this
chapter cites no number from them.

\paragraph{Verdict.} No index-calculus or generic measurement attributes a
cost difference to the conductor. Every apparent effect so far was
accounting (round 2's A0), sampling (the round-1 lead), the coordinate
model (Table~\ref{tab:iso-volcano-v1} and \S\ref{sec:corr-structure}) or
the automorphism group of a $j=0$ or $j=1728$ crater.

\section{Isogeny-derived factor bases}\label{sec:iso-fb}

By Lemma~\ref{lem:conservation}, a factor base built from an isogeny is a
base on one curve with a particular point set and presentation. Three
studies tested such bases against a non-isogeny control.
\begin{itemize}
\item \emph{Horizontal isogenies at $p=2^{61}-1$.} On a prime-order curve
  over $\FF_{2^{61}-1}$ (cofactor 1, trace $-1604388101$), all 15 outgoing
  kernel-distinct rational isogenies of prime degree at most 61 were
  built, one hop only, and the factor base was mapped through each
  (\art{research/prime61_isogeny_20261016/README.md}). Vertical separable
  isogenies do not exist in this class ($f_\pi=f_E=1$). All 16 workloads,
  12 paired streams each, reached their caps without a relation: a null
  relation-stage result, not a measured improvement. Coverage of the whole
  horizontal class is blocked on memory.
\item \emph{Isogeny kernels against Petit's $\mu_k$.} At $m=2$ and toy
  primes, the isogeny-kernel base matches a random base in yield and costs
  more only through its presentation (\S\ref{sec:pa-pkm}).
\item \emph{The PKM isogeny tower.} Its rows and their distance to rho are
  in Chapter~\ref{sec:prime} and the frontier table
  (Table~\ref{tab:ladder}); the seed of that tower is the next section.
\end{itemize}

\section{Isogeny seeds}\label{sec:iso-seed}

The algebraic arm \texttt{pkm-iso-t6-g4} builds its base as
$V=x(R'+\langle Q\rangle)$ on an auxiliary curve, and the seed selects the
auxiliary curve and the coset; the F4 system has the same shape for every
seed. A 3 November study selected the seed by a log-free coverage count $J$
on Curve25519-family miniatures
(\art{research/curve25519_isogeny_seed_20261103/README.md}).
\begin{itemize}
\item \emph{Census} (exact, 2{,}000 seeds per cell, $t=4$--$7$): the
  selected seed covers $1.62$--$3.24$ times the mean seed in all 15 cells
  with $t\ge5$ and equals the oracle-best seed in all 20 cells. At $k=20$,
  $t=6$, seed 1093 covers 528 points against seed~0's 172 ($3.07\times$);
  seed~0 sits at the 6th percentile.
\item \emph{Engine}: seed 1093 needs $1.75\cdot10^8$ F4 multiply-adds per
  relation against $5.42\cdot10^8$, $3.10\times$ fewer, at the same degree
  (5) and per-target cost ($0.991\times$).
\item \emph{End to end} (ecbench session \texttt{ECBS1h405f20219036}, 6
  targets, 48 of 48 solves verified): $S=26.448$ against $29.892$ at $w=0$,
  and $1.073\cdot10^6$ against $1.674\cdot10^6$ at $w=0.1$; negation rho
  is at $1.447$, $1.547$ and $2.448$ across the band and the projected
  generic baseline at $12.053$, $12.931$ and $20.831$. The paired ratio is
  $0.885$ $[0.570,1.373]$ at $w=0$ and $0.641$ $[0.379,1.084]$ at
  $w=0.1$ and $1$.
\end{itemize}
The preregistered criterion V2 (a mean ratio at most $0.6$) failed, and the
interval includes 1, so the session alone does not establish the gain; the
source reports that a later run (T9) replicated it on 16 targets. Two
details matter for the reading. Seed 1093 has 46 columns against 25, so it
needs about $1.84\times$ more relations, which is why $3.1\times$ cheaper
relations net about $1.7\times$; and the seed search, which ecbench does not
see, cost $1.35\cdot10^7$ group operations, which charged to one target
makes the band order unresolved and amortised per curve leaves it ordered.
Neither version changes the order against rho: the algebraic arm with seed 1093 is
$18\times$ negation rho at $w=0$ and about $4\cdot10^6\times$ at $w=1$. The corrections chapter records the
failed V2 (Table~\ref{tab:corr-transfer}).

\section{Isogeny computation as tooling}\label{sec:iso-tooling}

Several studies built isogeny machinery without an ECDLP claim; they are
recorded because later studies rest on them.
\begin{itemize}
\item \emph{Large prime degrees.} Every algorithm of Ben\v{c}ina, Kutas,
  Merz, Petit, Stopar and Weitk\"amper on fixed-degree isogenies between
  supersingular curves, and all of $\sqrt{\text{\'elu}}$ (Bernstein, De Feo,
  Leroux and Smith) beyond the codomain-and-map core that existed before,
  were implemented end to end and tested against an independent
  computation, except the quantum speed-ups, for which there is only a
  query-count model (\art{research/fixed_degree_isogenies_20261016/README.md}).
  The $\sqrt{\text{\'elu}}$ kernel-cost exponent $0.822$ $[0.805,0.838]$ is in
  Chapter~\ref{sec:transfer}.
\item \emph{Modular polynomials.} CM modular polynomials were built by
  explicit CRT (CM order selection, norm-equation primes, certified Hilbert
  polynomial seeds, paired class actions on surface and floor), with an
  exact-equality acceptance run; no speed is claimed
  (\art{research/cm_volcano_crt_20261016/REPORT.md}). The $n=83$ computation of
  $\Phi_{6473}$ in \S\ref{sec:iso-volcano} is the largest use.
\item \emph{An isogenous Boolean system.} Over $\FF_{2^{37}}$ a separable
  degree-73 isogeny of a Koblitz curve (74 rational cyclic kernels, a
  squarefree degree-36 kernel polynomial) gave an isogenous neighbour whose
  Boolean decomposition systems were compared with the source's: all 128
  systems completed with every planted decomposition verified on both
  curves, F5 matrix work fell by at least 10\% in 4 of 8 registered cells and
  rose by up to $2.00\times$ in another, and the preregistered criterion failed
  (\art{research/koblitz_isogeny_descent_37_results_20260925/RESULTS.md}). By
  Lemma~\ref{lem:conservation} the two systems describe the same relations;
  a presentation can change the solver's work in either direction.
\item \emph{Isogeny-loop endomorphisms.} Round~2 of the isogeny-loop GLV
  work corrected round~1's pricing of the Verschiebung in characteristic
  two, $V=[2]\circ F^{-1}$ (a doubling composed with square roots, not two
  squarings), which turns round~1's lead into a provable loss
  (\art{research/isogeny_loop_glv_20261011/README.md};
  Table~\ref{tab:corr-rho}); the four-dimensional GLV--GLS speed-up of
  $1.09$--$1.15\times$ is in Chapter~\ref{sec:transfer}. These speed up curve
  arithmetic, not the discrete logarithm.
\end{itemize}

\section{What the isogeny side settles}\label{sec:iso-summary}

\begin{itemize}
\item On prime fields the isogeny class contains no weak curve, by Tate;
  walks of standard curves are tooling.
\item Generic cost is level-blind to within a few per cent on tall
  volcanoes; the only level effect is the crater's automorphism group.
\item Index-calculus yield follows the factor base's point set, which
  depends on the coordinate model more than on the level; no conductor
  effect has survived its controls.
\item Isogeny-built factor bases are single-curve bases with a
  presentation cost; selecting the seed is a constant-factor lever that
  leaves the algebraic arm orders of magnitude above rho.
\end{itemize}
\chapter{Twists, pairings and lifts}\label{ch:twistspairings}

Three classical transfers move the discrete logarithm out of the curve
group: to the quadratic twist's eigenspace of $E(\FF_{q^2})$, to a
multiplicative group by a pairing, and to a $p$-adic or global lift.
Chapter~\ref{ch:obstructions} proves that twist points carry no relations
for the curve and records the lift-dense subspaces, the function-field
xedni bound and the oracle-simulation results. This chapter adds the
measured twist studies beyond that obstruction, the pairing and
embedding-degree closures, the anomalous-curve control, and the $p$-adic
and global-lift lanes with their corrections. Most items are theorem
closures of known results; they are recorded because each retired a
concrete proposal of the programme. Measured figures are producer-reported
unless a pcc claim is named, and no pcc standing was computed for this
draft.

\section{Twists beyond the obstruction}\label{sec:tp-twist}

Over $\FF_{q^2}$, $E(\FF_q)$ and the twist $E'(\FF_q)$ are the $+1$ and
$-1$ Frobenius eigenspaces of $E(\FF_{q^2})$, and $\#E(\FF_{q^2})=\#E(\FF_q)\cdot\#E'(\FF_q)$.
Applying Frobenius to a decomposition $R=A+B$ of $R\in E(\FF_q)$, with $A$
a sum of $E$-points and $B$ of twist points, gives $2B=O$
(Chapter~\ref{ch:obstructions}). Three studies measured what that leaves.

\paragraph{Exact yield on toy prime fields.} The 20 October study added
twist abscissas to an $x\in\FF_p$ factor base on $y^2=x^3+2x+3$ and solved
$\Sem3$ for every non-2-torsion target, at $p\in\{101,211,233,307\}$ with
24 abscissas of each kind (\art{research/twist_index_calculus_20261020/REPORT.md}).
The union has 1{,}176 unordered candidate pairs per target against 300,
$3.92$ times as many, and gives \textbf{0} extra width-two hits on every
field, against 569--574 hits from $E$ alone. The twist abscissas do produce
relations, 569--573 per field, but for targets on the twist. Each curve has
the 2-torsion point $(-1,0)$; 3--7 twist pairs per field sum to it, and
every resulting three-factor equality reduces to $R=F+(-1,0)$, a shorter
relation plus a known translation. The source notes that counting raw roots
of a higher summation polynomial without lifting and reducing them can
overstate useful yield for exactly this reason.

\paragraph{A case where the twist is usable.} For $p\equiv2\pmod3$ the
supersingular curve $E_0:y^2=x^3+1$ has a rational 3-isogeny to
$Y^2=X^3-27$, which a rational scaling maps to a nontrivial quadratic twist
of $E_0$. The composition carries a prime-order subgroup to the twist; it
was checked on every ordered pair of points at $p=101$ and $233$ (10{,}404
and 54{,}756 homomorphism checks) and on a transported scalar. Width-two hit
counts on transported targets then differ in both directions with the base
(for example 92 against 95 at $p=101$, $b=24$, and 12 against 15 at $p=233$,
$b=16$, where more hits cover fewer targets). The source reads this as
transport and local count differences, not a speed-up; these curves also
have a small embedding degree (\S\ref{sec:tp-pairing}). It is an instance of
Lemma~\ref{lem:conservation}: an explicit rational map moves the target,
and only the destination's base geometry and arithmetic can differ.

\paragraph{Seven twist ideas, preregistered.} The 27 October study froze
one preregistration per idea before data and had each result replicated
and attacked by independent agents (\art{research/twist_ecdlp_20261027/README.md};
sizes up to 40 bits, operation counts primary). No arm that used the twist
beat its matched baseline (Table~\ref{tab:tp-twist}).

\begin{table}[t]
\centering\small
\begin{tabularx}{\linewidth}{@{}l>{\raggedright}p{3.6cm}>{\raggedright}p{2.6cm}>{\raggedright\arraybackslash}X@{}}
\toprule
& Idea & Status & Key number \\
\midrule
E0 & separation lemma $2B=O$ & supported & 0 violations in $8{,}254+34{,}498$ classes \\
E1 & mixed-solution census ($\Sem3$, $x$-only) & structure supported; rate inconclusive from $2^{20}$ & 0 mismatches over $91{,}200+508{,}800$ targets; rate ratio $0.9965$ (12 bit), $0.9935$ (16 bit), $0.85$ $[0.29,1.47]$ (20 bit) \\
E2 & joint $E+E'$ table for twin targets & gain refuted & $1.000$--$1.062\times$ separate runs \\
E3 & $\FF_{q^2}$ eigen-decomposition & a loss & $53.3$--$2{,}812\times$ more work per $E$-relation \\
E4 & Koblitz $K_0\leftrightarrow K_1$, $\tau$-orbit base & no gain & identical 773{,}404 relations at $2.0246\times$ the solves \\
E5 & Kummer-line multiplicative walk & inconclusive & $4.6$--$6.2\times$ the baseline; falsifiers fired \\
E6 & genus-2 gluing of $E\times E'$ & construction only & $\#J(C)=\#E\cdot\#E'$ on 580 of 580 \\
\bottomrule
\end{tabularx}
\caption[]{The twist ideas of the 27 October study. The genus-2 arm's ``no
gain'' comes from a cost model and is not measured. E1, E2, E4 and E6 had
replication mismatches that were diagnosed as incomplete replication runs
and were not re-replicated. Source: \art{research/twist_ecdlp_20261027/README.md}.}
\label{tab:tp-twist}
\end{table}

The reference arm G1, a negation-map rho, solved every instance; the
Kummer-line walk cost $1.287$--$5.907\times$ it, with no row below 1. The
study's open items are the mixed-solution rate at $2^{20}$--$2^{24}$ and
$m\ge4$, factor bases over extension fields, a measured genus-2 Jacobian
arm, and pricing every arm across the band.

\paragraph{A uniform $S_3$ census of twist transfers.} A long-running
lane (\art{research/quadratic_twist_evidence_20261021/README.md}) executed
an explicit rational dual of a twist transfer at $\FF_{13}$, a uniform
census (96 arms, 896 source recoveries, 210 policy comparisons, 0 of 210
strict joint gains in rank or degree) and complete size-two and size-three
studies at $\FF_{37}$: 336 arms at caps 0--12, 4{,}368 profiles and zero
strongest-source joint records, on eighteen repeated original points. Its
own status line keeps the original goal, a useful-yield or degree advantage
from an actual twist, open; the registry still lacks 1{,}921 of its study
rows. The lane proves conditional algebra and finite transport, not a
yield or speed statement, and this treatise cites it only for that.
Two November proposals (a follow-up and a chord preflight, the latter
\texttt{PCC-20261105-7179e4}) have not run.

\paragraph{Lift-dense lines, re-checked.} The obstruction chapter records
that on composite-degree binary fields suitable $\FF_{2^d}$-lines lift
entirely when $\operatorname{Tr}(a)=0$. Two 6 November studies checked
claims that arrived from another session as unverified assertions, each
with a frozen protocol and a fresh instrument:
\begin{itemize}
\item \emph{The 2-torsion pairing.} A valid line stable under adding the
  2-torsion point exists exactly when $\operatorname{Tr}_{N/d}(b^{1/4})=0$,
  and it is then $b^{1/4}\FF_q$ (48 of 90 toy curves, 0 mismatches); on
  those lines $x(P+T)=\sqrt b/x(P)$, so their columns pair and they give
  about $2\times$, not $2^n$. None of the valid lines of c2pnb176v1,
  208w1, 272w1, 304w1 and 368w1 (180, 180, 270, 210 and 330 lines,
  reproducing the October counts) is stable, so no earlier count changes
  (\art{research/lift_dense_torsion_pairing_20261106/README.md};
  \texttt{PCC-20261106-b5904a}, \texttt{PCC-20261106-74c86c}).
\item \emph{Odd characteristic.} For lines $x_0+u\FF_q$ on $y^2=f(x)$ over
  $\FF_{q^n}$, a line is all-lift or all-twist when, for every Frobenius
  orbit $O$ of the roots of the shifted cubic, $n\,r_O/s_O$ is even. All
  1{,}287 lines meeting this criterion are constant; at $(q,n)=(101,2)$
  and $(101,3)$ none of 3{,}520 lines failing it is; one chance-constant line
  occurred at $(13,3)$. On descent lines, which exist on subfield curves up
  to isomorphism, the dead fraction of lifted points is $0.505$--$0.615$ per
  cell, but the registered rule for \emph{which} half is dead mispredicts
  405 of 27{,}880 abscissae
  (\art{research/lift_dense_orbit_parity_20261106/README.md};
  \texttt{PCC-20261106-3e6fc4}).
\end{itemize}
These are the only items in this chapter with pcc claim records; the
producers label every statement in them unverified. They refine the count
of usable lines; they do not price a factor base against rho.

The order-24 twist action on $j=0$ curves over $\FF_{p^4}$, which changes
the matched rho on those curves, belongs to the extension-field programme
and is in \S\ref{sec:glv-order24}.

\section{Pairings and the embedding degree}\label{sec:tp-pairing}

\paragraph{MOV and Frey--R\"uck.} The Weil or Tate pairing with an
independent $N$-torsion point maps $\langle P\rangle$ injectively into
$\mu_N\subset\FF_{p^k}^*$ with $k=\operatorname{ord}_N(p)$, where the
discrete logarithm falls to finite-field index calculus. The atlas closes
the lead as known (\art{crypto-autoresearcher:knowledge/leads/ecdlp/transfer/LEAD-TRF-mov-small-embedding-degree.yaml}):
$k$ is a function of $(N,p)$ and so an isogeny-class invariant
(\S\ref{sec:iso-tate}), and for random prime-field curves it is of size
about $N$ (Balasubramanian--Koblitz). The transfer is a check on a curve,
not an ECDLP algorithm. The programme's scoping question for pairings,
\texttt{RQ-PAIR-21313f}, has toy checks specified and zero runs. Curve
constructions in this programme check the embedding degree as a matter of
course; the conductor-gap pairs of \S\ref{sec:iso-conductor}, for example,
all have embedding degree above 20.

\paragraph{Two routes to a second torsion direction.} Two proposals
tried to obtain MOV without a second independent point; both are closed by
theorem.
\begin{itemize}
\item \emph{A self-pairing from one point.} On a cyclic group of prime order
  every bilinear map is a multiple of one fixed map, so a nontrivial
  self-pairing on the rational $N$-torsion line needs all of $E[N]$
  rational, that is $N\mid p^k-1$ with $k=1$. An endomorphism of an
  ordinary curve commutes with Frobenius and preserves the rational
  eigenline, so it is not a distortion map
  (\art{crypto-autoresearcher:knowledge/leads/ecdlp/transfer/LEAD-TRF-one-point-self-pairing.yaml}).
\item \emph{Trace-zero varieties and $(N,N)$ graph kernels.} By
  Balasubramanian--Koblitz, $E(\FF_{p^k})[N]$ is larger than the rational
  line only if the embedding degree divides $k$; otherwise the trace-zero
  variety $T_k(E)[N]$ is zero, since trace acts as $[k]$ on the line. An
  $\FF_p$-rational $(N,N)$ kernel of the Kani--Robert type in $E\times E'$
  exists only when the embedding degree is 1. Both reduce to MOV
  (\art{crypto-autoresearcher:knowledge/leads/ecdlp/transfer/LEAD-TRF-second-torsion-direction.yaml}).
\end{itemize}

\paragraph{What stays open on the pairing side.} Two leads remain, neither
with a run.
\begin{itemize}
\item \emph{The target-field cost of pairing curves.} For BLS12-381 and
  BN254 the ECDLP is bounded by the DLP in $\FF_{p^{12}}$ after MOV; the
  quoted security levels come from one family of exTNFS estimates. A
  recomputation with memory charged has 33 experiment contracts and zero
  runs; its next step is a certified toy tower-field DLP and a same-shape
  calibration against a published record
  (\art{crypto-autoresearcher:knowledge/leads/ecdlp/transfer/LEAD-TRF-pairing-target-field-recost.yaml}).
  It bears only on pairing curves.
\item \emph{Miller-cocycle maps when the pairing is trivial.} With
  $\gcd(N,p-1)=1$ the Tate pairing is trivial, but the Miller function
  gives a canonical non-homomorphic map $U(R)=f(R)^{1/N}$. If $\log U(kP)$
  carried structure in $k$ it would be a cheap transfer into $\FF_p^*$;
  equidistribution at the $1/p$ rate would close it. The lead is
  unexamined
  (\art{crypto-autoresearcher:knowledge/leads/ecdlp/transfer/LEAD-TRF-miller-cocycle-nonhom-transfer.yaml}).
\end{itemize}

\section{Anomalous curves and Smart's attack}\label{sec:tp-smart}

When $\#E(\FF_p)=p$, lifting to $E(\mathbb{Q}_p)$ and applying the formal
logarithm makes the discrete logarithm linear in $\ZZ_p$ (Smart,
Satoh--Araki, Semaev). It works exactly on trace-one curves, an
isogeny-class invariant absent from generic and standard curves. The
programme used it as a positive control: in \texttt{EXP-IT-001} a Smart
certificate at $p=771853$ ($N=p$, trace 1) cost $C_{\rm special}=157=\lceil8\log_2p\rceil$
against a matched rho of 779, a ratio of $0.202$ below the control
threshold $0.7$ (\art{crypto-autoresearcher:ledger/evidence/EV-IT-aefd12.yaml}).
A toy study re-derived and replicated the known degeneracy: Smart's attack
fails exactly on the canonical-lift class modulo $p^2$, where the lift is
not random (\texttt{KN-FIND-3c7e19}, cited in
\art{crypto-autoresearcher:knowledge/leads/ecdlp/transfer/LEAD-TRF-anomalous-padic-lift.yaml}).
The atlas lists the lead as \texttt{closed\_known}, to be revisited only as
a positive control or for a curve with $\#E=p$.

\section{$p$-adic sections and the CZLIFT lane}\label{sec:tp-padic}

The natural generalisation of Smart to non-anomalous curves is a computable
section $E(\FF_p)\to E(\mathbb{Q}_p)$ whose formal logarithm or $p$-adic
digits linearise the logarithm.

\paragraph{The torsion section is the only homomorphic one.} For
$p\nmid\#E(\FF_p)$ the only homomorphic section is the prime-to-$p$
torsion section; the formal logarithm vanishes on torsion, as $\log_p$
kills the Teichm\"uller section of $\FF_p^*$. The programme proved this at
derivation tier, uniformly in $p\ge5$ (\texttt{KN-FIND-7e2b94}): the section
costs $O(\log n+k\log p)$ additions and Smart's functional vanishes on it,
so the ``full break'' clause of the hypothesis is false
(\art{crypto-autoresearcher:knowledge/leads/ecdlp/transfer/LEAD-TRF-padic-canonical-section.yaml}).
The toy instrument agreed on 1{,}116 of 1{,}116 and 576 of 576 pairs
(\texttt{EV-CZLIFT-ad28f5}). The one failure, a naive section of $\FF_p^*$,
was rejected with strength \emph{strong} and explained as a size effect
when $g^d<p$, not a $p$-adic one (\texttt{EV-CZLIFT-f2a401}; $p<2000$). Any
non-homomorphic section recovers $d$ only at the random $1/p$ rate. One
exponent-relevant question stays open: whether the maximal Fourier mass of
$p$-adic digit characters of the torsion section decays faster than the
$n^{-1/4}$ break-even (\texttt{KN-OPEN-d3a6f0}). The round synthesis
\texttt{DEC-20261024-23d64d} records that the question
\texttt{RQ-CZLIFT-b9f765} cannot close yet, with a control and a review
outstanding.

\paragraph{Two November follow-ups.} Two preregistered studies tested
specific lift-derived observables.
\begin{itemize}
\item \emph{H110, the $p^2$-lift defect.} A deterministic Hensel section
  into $E(\ZZ/p^2\ZZ)$ is executable, but the scalar polynomial of its
  first formal-layer defect had the maximum possible degree $r-2$ (81 to
  945) on all 28 scalar words, against frozen gates of 49 to 100, so 0 of
  28 passed; all 32 random-section controls and both coordinate-isomorphic
  controls were also at maximum degree. All 9{,}096 targets passed the
  arithmetic checks, and every solved torsion lift had $\tau=0$, which is
  the vanishing above (\art{research/p2_lift_defect_20261101/README.md}).
\item \emph{H109, a division-polynomial derivative oracle.} The invariant
  derivative of the order-$r$ division polynomial at the target was tested
  for a constant-size character formula for scalar parity. With
  10{,}674{,}040 field multiplications and 211{,}712 inversions charged and
  all 8{,}832 coefficient tuples enumerated, exact formulas were found on 0
  of 28 rows; the candidate's best-match rate was $0.500$--$0.611$ against
  $0.500$--$0.630$ for 32 symmetry-matched controls
  (\art{research/torsion_division_derivative_20261031/README.md}).
\end{itemize}
Each rejects only its frozen family; neither establishes a lower bound.

\section{Global lifts, heights and the CM canonical lift}\label{sec:tp-global}

Silverman's xedni calculus lifts $(E,P,Q)$ to a curve over $\mathbb{Q}$ or a
number field and hopes that the lifts satisfy a short Mordell--Weil
relation. The obstruction chapter gives the exact function-field lift
probability. Three items refine the global case.

\paragraph{The JKSST theorem, audited in Lean.} Jacobson, Koblitz,
Silverman, Stein and Teske's analysis derives the classic bound under
``certain plausible assumptions''. A kernel-checked audit of 1 November
separates what is proved from what is assumed
(\art{research/jksst_xedni_audit_20261101/README.md}). It repairs the
conditional theorem and refutes several overbroad readings:
\begin{itemize}
\item rank, radius and minimum separation do not determine a lattice ball
  count: two rank-two lattices with the same coarse parameters have
  radius-one counts 5 and 3 (sup-norm 9 and 3);
\item no fixed model discriminant bounds the N\'eron--Tate heights of
  \emph{all} multiples of a nontorsion point by one constant, which refutes
  only the universal reading, not the points xedni selects;
\item quadraticity, positivity and finite rank alone do not imply the
  uniform gap of Lang's height conjecture, which remains open.
\end{itemize}
The algorithm-specific bridges (selected-point heights against a global
minimal discriminant, Lang's conjecture, the random-image model, the
good-reduction correspondence) remain open; the audit does not claim the
opposite unconditional statement.

\paragraph{Measured toy lifts.} In the CZLIFT round of 23 October, small
rational points grew as predicted but did not reduce uniformly, and a
relation search found one relation, the tangent relation at $p=101$, in
3{,}200 valid trials with coefficients bounded by 6, and none of 1{,}600 at
rank 3; a planted dependence was recovered
(\texttt{EV-CZLIFT-ee9b23}). Three hypotheses were weakened and two run sets
declared inadmissible
(\art{crypto-autoresearcher:knowledge/leads/ecdlp/transfer/LEAD-TRF-global-lift-small-height.yaml}).
The lead stays open with a named next step: an explicit prime-field lift
family that either composes every bridge or falsifies one.

\paragraph{The CM canonical lift: a corrected closure.} Lane L9 asked
whether the complex multiplication of the Serre--Tate canonical lift shrinks
lifted heights enough for xedni. Its closure first rested on a norm-form
gap: $N(\alpha)\ge\abs{D_E}/4$ composed with
$\hat h(\alpha P)=N(\alpha)\hat h(P)$ would place every CM companion outside
the height budget. An independent red-team review found that this was not
the load-bearing argument, and correction \texttt{CORR-20260810-930c42}
(with \texttt{DEC-20260810-3110d5}) restated it
(\art{crypto-autoresearcher:knowledge/leads/ecdlp/transfer/LEAD-TRF-cm-canonical-lift-heights.yaml}).
The corrected closure covers ordinary prime-field Serre--Tate and global CM
models of fixed arity 2 to 9 with an absolute coefficient bound, and
twelve enumerated mechanisms. Its argument is a counting one: a
target-independent route chooses a $p^{o(1)}$ output set $S$ before $Q$, so
$\Pr[Q\in\operatorname{red}(S)]\le\abs S/N=p^{-1+o(1)}$, an instance-count
exponent of 1 against rho's $1/2$. For generic $\abs D$, writing the
canonical lift through the class polynomial costs degree $\abs D^{1/2+o(1)}$,
a separate rho-scale gate. The corrected scope excludes target-dependent
sections, growing arity or coefficient bounds, and the universal norm-form
claim, which was false. The conclusion survives; its reason changed.

\section{What the classical transfers settle}\label{sec:tp-summary}

\begin{itemize}
\item The twist contributes only 2-torsion to a decomposition, exactly;
  every measured twist arm is at or above its baseline, and the lift-dense
  refinements change counts of usable lines, not distances to rho.
\item Pairings help only at small embedding degree, which is a property of
  $(N,p)$; the post-MOV cost of deployed pairing curves is unmeasured here.
\item Smart's attack is a positive control. Every non-anomalous $p$-adic
  section the programme tried is either the torsion section, which the
  formal logarithm kills, or behaves like a random map.
\item The global-lift lane is held by conditional theorems with named open
  premises, one of them repaired in Lean and one corrected in scope.
\end{itemize}
\chapter{Weil descent and GHS covers}\label{ch:ghs}

The Gaudry--Hess--Smart construction~\cite{ghs2002} maps the discrete
logarithm on $E/\FF_{q^n}$ to the Jacobian of a curve over $\FF_q$ whose
genus is governed by the magic number $m$. Where the genus is small, index
calculus on the Jacobian is cheaper than rho on $E$. The treatise already
records what this means for ECC2K-130 (magic numbers 1 or 131, with the
Maurer--Menezes--Teske correction; \S\ref{sec:ecc2k130}), the measured
trial cost on a genus-255 Jacobian that turned the $N=130$ counterfactual's
margin negative (\S\ref{sec:ladder3}), the cubic GHS route to genus 3 and
its plane-quartic branch (Chapter~\ref{sec:transfer},
Table~\ref{tab:prog_transfer-quartic}), the censuses at $n=3$ and $n=5$,
and the Oakley groups (Chapter~\ref{sec:transfer}). This chapter adds three
bodies of work: the descent screens of the curve registry and the
corrections they went through, the exhaustive binary cubic catalog
campaign over $\FF_4$, $\FF_8$ and $\FF_{16}$, and the genus floors in odd
characteristic. Every attack cost in it is producer-reported, and most are
Magma wall times at toy field sizes.

\section{Descent screens on the registry}\label{sec:ghs-screens}

\paragraph{The scan.} A native screen read the registry of 321 curves (176
prime-field, 136 binary of which 72 random, 60 Koblitz and 4 subfield, and 9
extension-field toys) and enumerated 1{,}156 transfer routes, pricing each
with one frozen cost model (\art{research/jacobian_transfer_scan_20261010/README.md}).
Of the 321 curves, 314 were \texttt{boundary}, 3 \texttt{unresolved}, 2
\texttt{formula\_advantage} and 2 \texttt{formula\_marginal}. Every row is a
formula, not an attack:
\begin{itemize}
\item every prime-field curve, including the 54 standardised ones, had best
  ratio $-2.1$ to $-2.0$ bits, on the same-field genus-2 cover, whose
  Jacobian rho has rho's exponent and four times its operation cost;
\item every Koblitz, subfield and prime-degree binary curve was a boundary:
  Koblitz curves degenerate to $m=1$, random-$b$ prime-degree curves have
  $m=n$, and the 60 Koblitz curves sit at $-3.3$ to $-3.2$ bits on the
  genus-3 same-field cover;
\item the two formula advantages were Oakley Group~4 (magic number 5, genus
  15 over $\FF_{2^{37}}$, $+10.9$ bits) and c2pnb368w1 ($m=4$, genus 7 over
  $\FF_{2^{92}}$, $+10.1$ bits); Oakley Group~3 was marginal at $+1.1$ bits.
\end{itemize}

\paragraph{The magic-number correction.} Three days later an audit found
that the trapdoor module had defined the magic number as the plain orbit
dimension of $\sqrt b$ and decided the genus type by orbit length; both are
wrong against Maurer, Menezes and Teske, whose magic number is
$\dim\operatorname{span}\{(1,\sigma^i(\sqrt b))\}$ and whose type~I is
$(x-1)\nmid\operatorname{Ord}_b$ (\art{docs/audits/GHS_MAGIC_NUMBER_MMT_CORRECTION_20261013.md}).
The scan's magic numbers were already MMT's; only its type rule was wrong.
On regeneration no magic number changed, 252 GHS rows changed genus, 16
rows changed class and four curve verdicts changed
(Table~\ref{tab:ghs-mmt}). Every $n=2$, $m=2$ row moved from genus 1, not
carried, to genus 2, boundary. The ECC2K-130 consequences (the smallest
nontrivial genus over $\FF_2$ is $2^{130}-1$, not $2^{129}$) are in
\S\ref{sec:ecc2k130}.

\begin{table}[t]
\centering\small
\begin{tabular}{@{}lllll@{}}
\toprule
Curve & Best route & Genus & Ratio (bits) & Verdict \\
\midrule
Oakley Group 3 & GHS to $\FF_{2^{31}}$ & $15\to16$ & $+1.15\to-3.04$ & marginal $\to$ boundary \\
Oakley Group 4 & GHS to $\FF_{2^{37}}$ & $15\to16$ & $+10.94\to+6.76$ & advantage $\to$ marginal \\
c2pnb368w1 & isogeny-extension to $\FF_{2^{46}}$ & $7\to8$ & $+10.05\to+7.92$ & advantage $\to$ marginal \\
c2pnb272w1 & GHS to $\FF_{2^2}$, $m=8$ & $127\to128$ & $+41.32\to+72.20$ & unresolved $\to$ advantage \\
\bottomrule
\end{tabular}
\caption[]{The four curve verdicts changed by the GHS type correction of
13 October. Ratio: $\log_2(\text{rho})-\log_2(\text{route})$ under the
frozen cost model, positive when the formula prices the route below rho.
None is an attack. Source:
\art{docs/audits/GHS_MAGIC_NUMBER_MMT_CORRECTION_20261013.md} \S2.}
\label{tab:ghs-mmt}
\end{table}

A second commit replaced the screen's bracket for $\operatorname{Tr}(a)=1$
with even relative degree (``$m$ or $m+1$'') by MMT's Lemma~6 and its
condition~(8): $a$ never changes the magic number, but it decides whether
the descent exists (\art{research/jacobian_transfer_lemma6_20261013/README.md}).
Across 2{,}091 registry curves, three rows changed (c2pnb176w1, c2pnb272w1
and c2pnb304w1, now exact at $m=8$, genus 128); two of them changed
verdict, c2pnb176w1 to \texttt{formula\_advantage} ($+23.64$ bits) and
c2pnb304w1 to \texttt{formula\_marginal}. The c2pnb rows over $\FF_4$ are
the kind of formula row whose genus (128) is beyond any implemented
Jacobian arithmetic; their advantage is a model output.

\paragraph{The $N=130$ counterfactual, priced at the corrected genus.}
Over $\FF_{2^{130}}/\FF_2$ the old table priced genus 128 ($m=8$) at
$2^{36.41}$ against rho's $2^{60.31}$. Under MMT no $b$ has $m=8$; the
attainable values start $1,2,5,6,9,10,13,14,\dots$
(\art{research/ghs_counterfactual_n130_20261013/README.md}). The least
admissible outcome is $m=9$, type~I, genus 255
($\operatorname{Ord}_b=\Phi_5(x)^2$), and a concrete $b$ with these
invariants and condition~(8) was constructed. Its modelled cost is
$2^{55.60}$ (smoothness bound 27, relations $2^{54.46}$, linear algebra
$2^{54.60}$), a margin of $-4.71$ bits instead of $-23.91$, checked against
the legacy pricer to 0.01 bit. The model crosses rho between genus 288
($2^{59.91}$) and 296 ($2^{61.34}$). Table~\ref{tab:ghs-n130} gives every row.

\begin{table}[t]
\centering\small
\begin{tabular}{@{}rrlll@{}}
\toprule
$l$ & $n$ & old $(m,g)$, cost & now $(m,\text{type},g)$, cost & below rho now \\
\midrule
1 & 130 & $(8,128)$, $2^{36.41}$ & $(9,\mathrm{I},255)$, $2^{55.60}$ & yes (model) \\
2 & 65 & $(7,64)$, $2^{299.93}$ & $(13,\mathrm{I},4095)$, $2^{43242.05}$ & no \\
13 & 10 & $(5,16)$, $2^{68.63}$ & $(5,\mathrm{I},15)$, $2^{64.52}$ & no \\
26 & 5 & $(4,8)$, $2^{60.80}$ & $(5,\mathrm{I},15)$, $2^{88.78}$ & no \\
\bottomrule
\end{tabular}
\caption[]{The $N=130$ GHS counterfactual recomputed under the MMT magic
number, against rho's $2^{60.31}$. Rows with $l\ge2$ use a fixed-genus model
the source calls pessimistic for large genus, so ``no'' means ``not shown to
break''. The $l=5$, 10 and 65 rows, absent from the old table, are not below
rho either. Source: \art{research/ghs_counterfactual_n130_20261013/README.md}.}
\label{tab:ghs-n130}
\end{table} A second, independent error in the
old table took the attainable magic numbers for $l\ge2$ from the
factorisation of $x^n-1$ over $\FF_{2^l}$ instead of $\FF_2$; it changes the
$l=2$ and $l=26$ rows, which are not below rho either way. The source's
own interpretation is that the remaining 4.71-bit margin does not absorb
its estimated $2^7$ curve additions per relation trial; measuring that trial
moved the margin to $-1.83$ bits (\S\ref{sec:ladder3}).

\paragraph{Other screens.} Three further studies bound what the binary
screens can claim.
\begin{itemize}
\item \emph{Generalised GHS.} Explicit maps for Hess's generalised binary
  construction at degree 12 gave 147 screened profiles of genus at most two
  and 232 index-two genus-two profiles, all members of known families;
  matched negation rho was run on each source, and the study states that
  the measurements do not support a calibrated GHS advantage over rho
  (\art{research/ghs_generalized_20261024/REPORT.md}).
\item \emph{A first comparison.} The October 9 GHS round had no higher-genus
  solver and so no measured GHS cost; its rho against BSGS raw-cost ratios
  were $2.35$ and $4.13$ (\art{research/ghs_curve_comparison_20261009/RESULT.md}).
\item \emph{Teske-type hidden descents.} A trapdoor in Teske's sense needs a
  weakness that is not a class invariant, a transfer along the isogeny and
  a weak set sparse in the class
  (\art{research/teske_variants_field_types_20261018/REPORT.md}). Prime
  fields have none of the three. Over binary fields the sparse regime is at
  relative degree $n=7$ for the type-I genus-7 descent only (margin $+0.5$,
  $+11.5$ bits at $N=161$), while the type-II genus-8 descent and Hess's
  variant are dense; at $n=15$, $21$, $31$ the public sub-descents (for
  example $n=5$ inside $n=15$) already beat rho. In odd characteristic the
  analogue of Teske's field is $\FF_{q^5}$ with a genus-5 family, by a
  count of weak sets of $\approx q^{r-3}$ curves reproduced at $q\le13$.
\end{itemize}

\section{The binary cubic GHS catalog campaign}\label{sec:ghs-catalog}

From 26 to 28 October about 25 studies classified, exhaustively, which
ordinary elliptic curves over $K=\FF_{q^3}$ with $q\in\{4,8,16\}$ admit a
genus-3 or genus-4 Jacobian transfer to $k=\FF_q$, constructed explicit maps
where they exist, and ran complete attacks against a matched rho on the
source subgroup for a few. The treatise cites only the genus-3 subcover
study among them (Table~\ref{tab:distance-weak}). The campaign is a closed
unit: an exact classification with uniformly negative toy attacks.

\paragraph{The census over $\FF_{16}$.} For $k=\FF_{16}$ and $K=\FF_{4096}$
the normalised ordinary cohort has 8{,}190 $K$-isomorphism classes and 128
group orders (\art{research/ghs_cubic_catalog_q16_20261026/REPORT.md}). The
LMFDB catalogue of abelian varieties records a genus-3 Jacobian in 126 of
the 128 cubic Weil-restriction isogeny classes, which contain 8{,}142 of
the sources. Of these, 5{,}534 have no direct binary cubic genus-3 GHS cover
even after arbitrary $K$-isogeny walks, and 5{,}302 of those pass a
simple-root gate for an isogeny nonzero on the largest prime subgroup. The
source calls this a catalog-backed eligibility bound, not an attack: the
Jacobians are already catalogued and Tate's isogeny is not evaluated.
Among the catalog-positive models with order $\equiv4$ or $6\pmod8$, outside
the Menezes--Teske walk set defined by order $\equiv0$ or $2\pmod 8$, are
4{,}032 class-C models, 3{,}852 of which pass the gate; this was noticed after
the frozen join and recomputed in the final audit.

\begin{table}[t]
\centering\small
\begin{tabular}{@{}lrr@{}}
\toprule
Class of source & Orders & Models \\
\midrule
all normalised ordinary curves over $\FF_{4096}$ & 128 & 8{,}190 \\
catalogued genus-3 Jacobian in the Weil-restriction class & 126 & 8{,}142 \\
\quad no direct cubic genus-3 GHS cover, after any $K$-isogeny walk & & 5{,}534 \\
\quad and the largest prime subgroup passes the simple-root gate & & 5{,}302 \\
\addlinespace
irreducible Frobenius polynomial (simple restriction) & 120 & 7{,}584 \\
\quad outside the Hess genus-3 closure & 46 & 2{,}691 \\
\quad\quad with a catalogued genus-3 Jacobian & 44 & 2{,}643 \\
\quad\quad without one & 2 & 48 \\
reducible Frobenius polynomial & 8 & 606 \\
\addlinespace
excluded for all of $\ell=3,5,7,11,13$ (one-step $(\ell,\ell,\ell)$) & 2 & 72 \\
\bottomrule
\end{tabular}
\caption[]{The exact partition of the $\FF_{16}$ cubic census. The first
block is the catalog join, the second the simple-restriction bound with the
Hess closure, the third the odd-isogeny screen on the 44 catalogued simple
residual classes. Sources under \art{research/}:
\texttt{ghs\_cubic\_catalog\_q16\_20261026/},
\texttt{ghs\_q16\_simple\_restriction\_bound\_20261027/},
\texttt{ghs\_q16\_odd\_isogeny\_screen\_20261027/}.}
\label{tab:ghs-partition}
\end{table}

\paragraph{What the structure allows.} Four exact results narrow the
search (Table~\ref{tab:ghs-partition}).
\begin{itemize}
\item \emph{Hess's construction, closed.} Within Hess's Artin--Schreier
  representation the only genus-3 Frobenius-order pairs are $(1,2)$ and
  $(2,2)$, which gives an if-and-only-if criterion: $\delta=\sqrt b\in P$
  when $\operatorname{Tr}(a)=0$ and $\delta\in U^*$ when
  $\operatorname{Tr}(a)=1$, with $\abs{U^*}=q^2-1=255$ and
  $\abs P=(q-1)(q+1)(q+2)/2=2{,}295$
  (\art{research/ghs_cubic_hess_closure_q16_20261027/REPORT.md}).
\item \emph{A genus lower bound for any correspondence.} If the Frobenius
  polynomial $T^6-tT^3+4096$ of $\operatorname{Res}_{K/k}E$ is irreducible,
  the restriction is simple and any nonzero homomorphism to a Jacobian
  needs genus at least 3; irreducibility holds for 120 of the 128 orders
  (7{,}584 models). Independently, a prime-order target $r>625$ needs genus
  at least 3 by the Weil bound $r\le25^g$
  (\art{research/ghs_q16_simple_restriction_bound_20261027/REPORT.md}).
\item \emph{One-step odd isogenies.} A rational $(\ell,\ell,\ell)$ isogeny
  needs a monic cubic divisor of the Frobenius polynomial modulo $\ell$.
  Among the 44 simple classes outside the Hess closure with a catalogued
  genus-3 Jacobian, $\ell=3$ excludes 18 classes and $\ell=5$, 7, 11, 13
  exclude 21--23; exactly two classes (orders 4076 and 4118, 72 models) fail
  for all five primes. At $\ell=3$ there is a closed form: for $q=4^m$, a
  trace $t\equiv0\pmod3$ excludes every rational one-step $(3,3,3)$ isogeny
  (\art{research/ghs_q16_odd_isogeny_screen_20261027/REPORT.md}).
\item \emph{Quotients and neighbours.} The four hyperelliptic genus-3 curves
  in the class of the order-3974 source have no separable degree-2 map to an
  elliptic curve (\art{research/ghs_cubic_quotient_obstruction_20261026/REPORT.md});
  automorphism and degree-3 neighbour screens at $q=4$, 8 and 16 exclude
  further direct covers, including the order-526 and 542 tails at $q=8$.
\end{itemize}
A scalar-3 Frobenius gate for a $(3,3,3)$ transfer is self-labelled a
rediscovery: the repository's prime-characteristic cubic-transfer work had
already derived the scalar-3 and Jordan module screen; the binary study
extends it to $q=16$ with an exact count for every $q=4^m$
(\art{research/ghs_scalar3_frobenius_gate_20261027/REPORT.md}).

\paragraph{Witnesses.} Searches produced explicit curves in the classes
that matter.
\begin{itemize}
\item Over $\FF_4$, 2{,}048 drawn plane quartics gave 55 smooth ones with
  the right point counts, 48 of which match an elliptic source over
  $\FF_{64}$, in 12 order classes; 8 of the 12 have no direct genus-3
  source. The LMFDB already lists Jacobians, including plane quartics, in
  all eight classes, so no new eligible class follows
  (\art{research/ghs_binary_cubic_isogeny_20261026/REPORT.md}).
\item Over $\FF_{16}$, three class-C source orders gained plane-quartic
  witnesses, and a 30{,}000-draw search on 13 residual orders found four
  smooth quartics in matching classes (orders 4006, 4054, 4078, 4126); the
  source-to-quartic maps were not built
  (\art{research/ghs_q16_residual_quartic_search_20261027/REPORT.md}).
\item The cubic Tate bridge bounds the number of $\FF_{4096}$ classes for
  which a complete attack beats rho between 0 and 2{,}608 (orders in
  $[3969,4225]$): an eligibility count, not a measurement
  (\art{research/ghs_cubic_tate_bridge_20261026/REPORT.md}).
\item Four genus-4 Hess curves have explicit degree-4 maps to their
  elliptic complements, with Lean identities
  (\art{research/ghs_q16_hess_elliptic_quotient_20261027/REPORT.md}).
\end{itemize}

\paragraph{The end-to-end attacks.} Table~\ref{tab:ghs-e2e} lists every
complete transfer attack of the campaign against a rho on the same source
subgroup and target. Every one loses. The range is wider than ``216 to
2{,}784 times'': one genus-4 representative reads $198\times$, the
order-3978 representative, with an order-17 subgroup, $3{,}300\times$, and
the $\FF_8$ quartic panel about $623{,}000\times$.

\begin{table}[t]
\centering\small
\begin{tabularx}{\linewidth}{@{}>{\raggedright}p{4.1cm}>{\raggedright}X>{\raggedright\arraybackslash}p{2.6cm}@{}}
\toprule
Study & Instance & Attack / rho \\
\midrule
Hess genus 4, four residual sources & orders 4006, 4054, 4078, 4126; $r=2003$--$2063$; 512 rho trials each & $512$, $698$, $698$, $688\times$ \\
Genus-4 frontier & order 3974, $r=1987$: 13 relations, $0.22$\,s against $0.00102$\,s & $216\times$ \\
 & order 4178, $r=2089$ & $198\times$ \\
 & order 3978, $r=17$ & $3{,}300\times$ \\
Nonsplit, order 4222 & plane-quartic line sieve, same Magma backend & about $271\times$ \\
Nonsplit $V_4$, order 4054 & attack core $0.360$\,s CPU against $0.0016$\,s per rho solve & about $225\times$ \\
Genus-3 subcovers & 15 exact equations, orders 3970, 4098, 4162 (Sage) & $549$--$2{,}039\times$ \\
 & two nonzero-trace orbits, class \texttt{3.16.a\_a\_ah} & $1{,}236$--$1{,}430\times$ \\
Hess $c$-parameter audit, $q=8$ & $\#E=548$, $r=137$; $\#E=556$, $r=139$ & $2{,}304$, $2{,}784\times$ \\
$\FF_8$ quartic panel & Frobenius-cyclic quartic, rho median $0.000554$\,s & about $623{,}000\times$ \\
\bottomrule
\end{tabularx}
\caption[]{Complete GHS or cover attacks of the binary cubic campaign
against a rho on the same source subgroup and target, in Magma or Sage CPU
or wall time at $q\le16$, on unisolated hosts. Ratios of medians or means as each
source states them. Sources under \art{research/}:
\texttt{ghs\_q16\_4126\_hess\_g4\_20261027/},
\texttt{ghs\_q16\_genus4\_frontier\_20261028/},
\texttt{ghs\_nonsplit\_order4222\_20261028/},
\texttt{ghs\_nonsplit\_v4\_cubic\_20261027/},
\texttt{ghs\_q16\_genus3\_subcovers\_20261028/},
\texttt{ghs\_hess\_c\_parameter\_audit\_20261026/},
\texttt{ghs\_q8\_quartic\_cover\_panel\_20261026/}.}
\label{tab:ghs-e2e}
\end{table}

Three caveats apply to the whole table. The times are Magma
\texttt{Cputime} or \texttt{Realtime}, or Sage wall time, with hundredth-second granularity
against sub-millisecond rho solves, so a ratio carries a large resolution
error at the small end. The subgroups are of order about $2^{11}$, where
the constant terms of a cover pipeline (map construction, factor-base
set-up) dominate. And the rho is a plain or Floyd walk in the same
calculator, not the strongest reference. None of these favours rho; they
make the ratios imprecise, not reversible. On the asymptotic side the
genus-3 route is the one with an exponent below rho's ($q^{4/3}$ against
$q^{3/2}$), and its measured behaviour at larger $q$ is in
Chapter~\ref{sec:transfer}; genus 4 has exponent $q^{2-2/4}=q^{3/2}$, the
same as rho's on $E(\FF_{q^3})$, so only a constant could separate them.

\section{Odd characteristic: genus floors}\label{sec:ghs-odd}

\paragraph{Which degrees reach genus 3.} For the GHS compositum attached to
a degree-2 function $x$ on $E/\FF_{q^n}$, the genus is
$g=1+2^{m-2}(r-4)$, where $m$ is the $\FF_2$-rank of the conjugate branch
sets and $r$ the size of their union; it depends on the Frobenius-orbit
type of the branch locus, never on $q$
(\art{research/odd_char_descent_genus3_20261014/REPORT.md}, Theorems 1--2).
Genus 3 occurs only for $(m,r)\in\{(2,6),(3,5)\}$, so for $n\ge4$ with
$\FF_q(j)=\FF_{q^n}$ every $x$ gives genus at least 5. For $n=3$ the
hyperelliptic genus-3 family
$y^2=e\,g(x)(x-\alpha)(x-\alpha^q)$, $\alpha\in\FF_{q^3}\setminus\FF_q$,
exists for every odd $q$ (Momose--Chao; Joux and Vitse~\cite{jouxvitse2012}
call it the only one). The tool checked Joux--Vitse's explicit model on 81
instances with $p\le31$: on every one the $L$-polynomial is that of
$\operatorname{Res}_{\FF_{p^3}/\FF_p}E$. Explicit genus-3 builds, a plane
quartic over $\FF_q$ and a hyperelliptic model with the transfer executed by
Cantor arithmetic, are in \art{research/ghs_cubic_genus3_build_20261012/} and
\art{research/ghs_cubic_hyperelliptic_genus3_20261013/}; neither solves a
logarithm.

\paragraph{The genus floor reprices the registry.} The same derivation
gives floors of 9, 49 and 41 at $k=6$, 10 and 14 for any GHS cover of
$E/\FF_{p^k}$ with $\FF_p(j)=\FF_{p^k}$, where every \texttt{cover\_descent}
row of the registry had been priced at genus $k$
(\art{research/cover_descent_genus_floor_20261015/README.md}). Priced at the
floor, 599 of the 846 rows that were \texttt{formula\_advantage} stay so
(576 at $k=6$, 23 at $k=14$), 41 become marginal and 206 boundary (all of
$k=10$, most of $k=14$); curve verdicts move from 848, 50 and 1{,}201 to 601,
67 and 1{,}431 (advantage, marginal, boundary). The floor is a lower bound
for GHS covers only, and no cover is exhibited, so the rows stay
\texttt{proposed}. The corrections chapter lists this repricing
(Table~\ref{tab:corr-transfer}).

\paragraph{The census at $n=4$.} Chapter~\ref{sec:transfer} and
\S\ref{sec:ladder3} quote the floor 13 for odd-order $E/\FF_{q^4}$. The
census gives the full picture by 2-torsion type, exhaustive at $q=3,5$
(\art{research/ghs_cover_census_n4_20261024/README.md}):
\begin{itemize}
\item odd $\#E(K)$: 13 if $j\notin\FF_{q^2}$, else 37;
\item $E(K)[2]=\ZZ/2$: 25 on every curve, whether $\#E\equiv2\pmod4$ or
  $4\mid\#E$;
\item $E(K)[2]=(\ZZ/2)^2$: 5 on most curves (26 of 28 at $q=3$, 158 of 210 at
  $q=5$), else 7 or 9.
\end{itemize}
Two preregistered predictions failed (floor 21 for $\ZZ/2$, observed 25;
odd-class minima in $\{13,25\}$, observed $\{13,37\}$), both from errors in
the producer's proofs found by the first run; after the amendment there were
0 violations over 1{,}256, 166 and 23 curve models at $q=5$, 3 and 7. A
structure counts as a descent to $\FF_q$ only if the four conjugate quotients
are distinct; $\sigma^2$-stable branch sets, with genus as low as 2, descend
$E$ only to $\FF_{q^2}$ and are excluded. Genus 5 at $n=4$ is $q^{8/5}$ in the
index-calculus exponent against rho's $q^2$, but only on curves with full
rational 2-torsion.

\paragraph{Decomposition coverage over $\FF_{q^n}$.} For a generic
$E/\FF_{q^n}$ and the base $x\in\FF_q$, the share of points that are a sum of
$n$ base points was computed exactly over the whole group
(\art{research/weil_descent_angles_20261018/REPORT.md}). In odd
characteristic at $n=2$--$5$ it matches a random $x$-set of the same size
(paired residual $\abs z\le0.94$). The exception is binary $n=2$, where the
coverage tends to $3/8$ instead of the Poisson value $1-e^{-1/2}=0.3935$
($z=-4.72$), for a proved reason: every $x\in\FF_q$ has trace zero over
$\FF_2$, so the base lies in one coset of $2E$ and only half the group can
decompose. For an attack the relevant figure is coverage of $2E$, which is
exactly $3/4$, about $1.9\times$ more relations per target than the model;
the source classes this as a cost-model correction. End to end at $n=2$ the
decomposition route still costs $2.7$--$66\times$ rho in every pricing
cell. The Kummer and Dickson tower bases on prime fields show no coverage
effect either (18 cells, $\abs z\le1.84$). A genus-2 route at $n=2$ (a
Scholten-type curve with $J(C)\sim\operatorname{Res}E$) needs about $p/3$
Jacobian operations against rho's $0.886\,p/\sqrt h$; its ratio to rho spans
$0.51$--$2.83$ across the band, so it is unresolved at every size, and at
40 bits rho is cheaper in 10 of 12 cells.

\paragraph{Cubic extensions in odd characteristic.} The native follow-up of
a historical Sage study certified the rational 3-torsion action and ran 58 native stage and
control runs; a complete native transfer is missing
(\art{research/cubic_extension_transfer_20261025/}; pcc claims
\texttt{PCC-20261105-a43948} and \texttt{PCC-20261105-a6723f}, a development
and a proposal, standing not computed here).

\section{Trapdoors in the isogeny class}\label{sec:ghs-teske}

Teske's trapdoor uses the fact that descent weakness is not a class
invariant over composite binary fields: a GHS-weak curve over
$\FF_{2^{161}}$ is hidden behind an isogeny that only its holder knows.
\S\ref{sec:ghs-screens} locates the fields where such a weakness can be sparse.
Two atlas leads ask the same question elsewhere.
\begin{itemize}
\item \emph{Prime fields.} By \S\ref{sec:iso-tate} no prime-field variant via
  descent exists; what remains is whether any non-generic index-calculus
  meter depends on the endomorphism ring
  (\art{crypto-autoresearcher:knowledge/leads/ecdlp/transfer/LEAD-TRF-teske-secret-isogeny-trapdoor.yaml}).
  A heavy-tail screen at 40-bit $N$ found five classes homogeneous
  (\texttt{EV-ECTD-8c098d}, which also discloses a fabrication defect in
  its rho-substitution prose); a vertical screen was invalid because its
  meter read the same vector on all 32 curves (\texttt{EV-ECTD-4c8a3b}); and
  every trial of the approved 20-arm panel \texttt{EXP-ECTD-0cf008} died with
  an infrastructure error. The isogeny-fibre pullback theorem, by which
  summation polynomials pull back along isogenies, predicts homogeneity,
  and the conductor-gap results of \S\ref{sec:iso-conductor} agree.
\item \emph{The Joux--Vitse weak class over $\FF_{p^6}$.} Membership in the
  weak class is not class-invariant (a random full-2-torsion curve lies in a
  weak class with probability $0.59$--$0.63$ at $p=11$ to $47$, \S\ref{sec:jac-jv}), so a hidden weak curve is structurally possible. The
  open question is whether some classes place their weak curve beyond a
  cheap public walk; the lead is unexamined
  (\art{crypto-autoresearcher:knowledge/leads/ecdlp/transfer/LEAD-TRF-teske-trapdoor-fp6-weak-class.yaml}).
\end{itemize}
A trapdoor of either kind would be a construction for a designer, not an
attack on a given curve.

\section{What descent settles}\label{sec:ghs-summary}

\begin{itemize}
\item On the registry only composite-degree binary curves have descent rows
  with a formula advantage, and every one is a model at a genus no
  implemented Jacobian arithmetic reaches; two type and admissibility
  corrections moved those rows in both directions.
\item Over $\FF_{4096}$ the classification is exact: which sources have a
  genus-3 carrier, which isogeny steps are excluded, and which classes
  remain. Every complete attack measured on it loses to rho, by factors from
  about 200 to about $6\cdot10^5$ at toy sizes.
\item In odd characteristic the genus is fixed by the branch-orbit type; it
  is 3 only at $n=3$, and the floors at $n=4$ depend on the 2-torsion, so
  prime-order-like curves over $\FF_{q^4}$ need genus 13 or 37.
\end{itemize}
\chapter{Jacobians of genus 2 to 5}\label{ch:jacobians}

Moving the problem into a Jacobian helps only when the Jacobian's index
calculus has a better exponent than rho on the source and the transfer is
cheap. Genus-$g$ double-large-prime index calculus~\cite{gtthd2007} costs
$\tilde O(q^{2-2/g})$ over $\FF_q$, and Diem's index calculus on a smooth
plane curve of degree $d$~\cite{diem2006,dt2008} costs $\tilde O(q^{2-2/(d-2)})$,
$\tilde O(q)$ for a plane quartic. For a source over $\FF_{q^n}$ the
Jacobian has dimension at least $n$, and rho costs $q^{n/2}$: at $n=3$ the
genus-3 routes have exponents $4/3$ and $1$ against $3/2$, at $n=4$ genus 4
and genus 5 have $3/2$ and $8/5$ against $2$. On a prime field the same
argument runs the other way. Chapter~\ref{sec:transfer} records the
genus-3 hyperelliptic route against the strongest rho, the plane-quartic
studies N4 and N4~v2 (Table~\ref{tab:prog_transfer-quartic}), Smith's walk,
and a summary of the Joux--Vitse cover; the rho references on Jacobians
are in \S\ref{sec:rv-jacobian}. This chapter adds the $n=3$ Tian route and
the N4 iteration between those two studies, the genus-4 and genus-5
Weil-restriction Jacobians, the same-field embeddings on prime fields, and
the full Joux--Vitse account, with its validated sieve row (\S\ref{sec:jac-jv}).

\section{Plane quartics at $n=3$}\label{sec:jac-quartic}

\paragraph{Theorem 4: when a Tian kernel exists.} Tian's construction
reaches a genus-3 Jacobian from $\operatorname{Res}_{\FF_{q^3}/\FF_q}E$ through
an isogeny with a Frobenius-stable kernel. The 18 October census proved when
such a kernel exists (\art{research/weil_restriction_isogeny_census_20261018/REPORT.md},
Theorem 4): for $\ell\ne p$ prime it exists if and only if
$\chi_\varphi(T)=T^2-tT+q^3$ has a root in $\FF_\ell$ and every root is a
cube in $\FF_\ell$, except that a Jordan $\varphi$ (a double root, $\varphi$ not
scalar) needs $\ell\equiv1\pmod3$ as well. So a nonsplit $\varphi$ never
qualifies, a split $\operatorname{diag}(\lambda,\mu)$ qualifies exactly when
$\lambda$ is a cube, and a scalar $\varphi$ always does. The census confirms
every cell; the split, scalar and nonsplit cases are proved for all $\ell$.

\paragraph{The Tian route measured.} A lead agent built the tools and ran
every measurement under a frozen protocol; no independent validator has
replayed the runs (\art{research/tian_n3_quartic_route_20261024/REPORT.md}).
\begin{itemize}
\item \emph{Census.} For every prime-order class at $q=7,11,13,17$ that
  Theorem 4 flags, the isogeny class of $\operatorname{Res}E$ contains the
  Jacobian of a smooth plane quartic over $\FF_q$: 26 of 26. Unflagged
  prime-order classes are hit too, at about half the rate, so a quartic in
  the class is not evidence of a Tian kernel. No isogeny was computed.
\item \emph{Diem index calculus on generic quartics.} A complete
  double-large-prime index calculus solved and verified every planted
  instance up to $q=8191$. With native work priced at 0 it costs 180--340
  group additions against rho's 5{,}000--300{,}000; at $w=0.1$ and $1$ it
  loses, because its dense linear algebra grows like $q^{1.4}$. The order
  does not hold across the band, so the ranking is unresolved.
\end{itemize}
The Diem arm ran on generic quartics, not on the flagged classes, and no
map from a source curve to its quartic was built. The benchmark copy
\art{benchmarks/20261024-tian-n3-quartic-ic/} has no validation file.

\paragraph{The step between N4 and N4~v2.} The 21 October iteration added
double large primes (B2) and structured elimination (B1) to the frozen N4
baseline B0 on the GHS-weak family $E_0/\FF_{q^3}$ with $4\mid\#E_0$
(\art{research/quartic_n4_la_20261021/README.md}). B2's lines through two
base points meet the quartic in two more rational points, which become
large-prime edges; the base shrinks from $3.0\,q^{0.6}$ to $2.0\,q^{0.5}$
points, and both phases scale as about $q^{1.0}$. All 57 B2 units verified.
The paired ratio B2/B0 is $0.355$ $[0.324,0.392]$ at $q=4051$, $0.224$
$[0.222,0.227]$ at 65{,}479 and $0.297$ $[0.295,0.299]$ at 262{,}111. Against a
matched \emph{plain} rho, IC/rho falls to $0.26$ $[0.23,0.29]$ at $q=4051$
and $0.024$ ($0.02$--$0.03$) at $q=1{,}048{,}507$ ($\ell\approx2.8\cdot10^{17}$),
resolved below rho across the band for every $q\ge991$; at $q=499$ the band
leaves it unresolved and at $q=223$ rho is lower. The fitted exponent on
$q\ge4051$ is $1.045$ $[1.037,1.052]$ against rho's $1.476$ $[1.438,1.508]$.
Table~\ref{tab:jac-n4la} gives the ladder. B1 is a trade ($0.81$--$0.91$ of B0 up to $q=65{,}479$, $1.53$ at 262{,}111) and
was not adopted. B2 became the baseline that N4~v2 then improved with a
negation-map rho; the host was a shared VM and there is no validation file
(\art{benchmarks/20261021-quartic-n4-la/}).

\begin{table}[t]
\centering\small
\begin{tabular}{@{}rrrrrl@{}}
\toprule
$q$ & units & rho runs & IC/rho, $w=0$ & $w=1$ & band \\
\midrule
223 & 5/5 & 150 & $3.85$ $[2.67,5.25]$ & $2.51$ & rho lower \\
499 & 5/5 & 150 & $1.40$ $[1.08,1.79]$ & $0.85$ & unresolved \\
991 & 5/5 & 150 & $0.79$ $[0.71,0.88]$ & $0.52$ & IC lower \\
4051 & 5/5 & 150 & $0.26$ $[0.23,0.29]$ & $0.18$ & IC lower \\
65479 & 5/5 & 80 & $0.05$ $[0.04,0.06]$ & $0.04$ & IC lower \\
262111 & 3/3 & 18 & $0.04$ $[0.03,0.05]$ & $0.02$ & IC lower \\
1048507 & 3/3 & 12 & $0.02$ $[0.02,0.03]$ & $0.02$ & IC lower \\
\bottomrule
\end{tabular}
\caption[]{N4 iteration 2, arm B2 (double large primes on lines through base
points), on the GHS-weak family $E_0/\FF_{q^3}$, $4\mid\#E_0$, against a matched
\emph{plain} rho, cold single-target cost in multiplication equivalents with
95\% bootstrap intervals; selected rows of thirteen. Band: the order across
$w=0$, $0.1$, $1$. Shared VM, no validation file. Source:
\art{research/quartic_n4_la_20261021/README.md}.}
\label{tab:jac-n4la}
\end{table}

\section{Weil restriction at $n=4$ and $n=5$}\label{sec:jac-weilres}

\paragraph{Genus-4 Jacobians.} For prime-order $E/\FF_{q^4}$ the Weil
polynomial of $\operatorname{Res}E$ is $T^8-tT^4+q^4$, and a curve matches
when its counts over $\FF_{q^k}$ for $k\le4$ do. A lead agent's census
(\art{research/genus4_n4_census_20261024/REPORT.md}) found that the class
contains a smooth non-hyperelliptic genus-4 Jacobian at every admissible
trace at $q=3$ (exhaustive, reproducing nine classes of an earlier round),
at 13 of 14 traces at $q=5$ (the exception is $t=49$) and at all 20 at $q=7$
(both sampled). No congruence on $(q\bmod\ell,t\bmod\ell)$ predicts which
classes are reached, and Tian-kernel existence does not either. Every curve
found has plane quintic models, the setting of Diem's plane-curve index
calculus at $\tilde O(q^{4/3})$ against rho's $q^2$, a cost the study cites
and does not verify. Its Theorem 6 is labelled an agent proof, checked
against direct linear algebra on 44 cells and not independently reviewed.
No map $J(C)\to E$ was computed: this is existence in the isogeny class, not
an attack.

\paragraph{The genus-5 cover.} A toy reproduction of the quartic-extension
genus-5 cover ran 90 complete attacks ($q=61$--$8009$, three seeds, full and
$\iota$-folded arms): all 90 logarithms recovered and every invariant
check passed, and matched rho was faster at every $q$
(\art{research/humbert_genus5_cover_20261021/README.md};
\art{benchmarks/20261026-humbert-genus5-cover/}, no validation file). A
first sweep contained one wrong logarithm at $q=1511$ from a linear-algebra
bug (unused factor-base columns) and is kept as superseded. The genus-5
exponent $8/5$ is below rho's $2$, so the loss at toy sizes is a constant,
and no crossover was measured or extrapolated.

\paragraph{Reading the two results together.} At $n=3$ the quartic route's
index calculus beats rho at $w=0$ on generic quartics but not across the
band; at $n=4$ and $5$ the Jacobians exist in the right classes, but either
no map is built (genus 4) or matched rho wins at every measured $q$
(genus 5). Neither touches a prime-order curve over a prime field.

\section{Same-field embeddings on prime fields}\label{sec:jac-samefield}

\paragraph{The closure.} For a prime-field source (Table~\ref{tab:jac-samefield}
summarises the measured collectors), a same-field genus-$g$
cover gives $\#J(C)(\FF_p)=N\cdot\#A(\FF_p)$ with $\#A=\Theta(N^{g-1})$, and
genus-$g$ index calculus costs $\tilde O(q^{2-2/g})\ge q$ for $g\ge2$: a full
exponent above rho's $N^{1/2}$. Quadric-pencil and quartic models of $E$
itself are isomorphic models with the same group. The atlas closes the
lead by theorem (\texttt{DEC-20260902-4c82fd};
\art{crypto-autoresearcher:knowledge/leads/ecdlp/transfer/LEAD-TRF-same-field-genus-g-embedding.yaml}),
and the registry scan's $-2.1$ to $-2.0$ bits on every prime-field curve is
its formula form (\S\ref{sec:ghs-screens}). The studies below measured
specific constructions against that closure.

\paragraph{Jacobians in ecbench.} The 25 October study added genus-2,
plane-quartic and GHS-cover families to ecbench
(\art{research/ecbench_jacobians_transfers_20261025/README.md}).
\begin{itemize}
\item On two genus-2 curves (17.5 and 20.2 bits), \texttt{ic.hyperelliptic}
  costs $1.473$ $[1.158,1.874]$ of matched negation rho, above rho at every
  weight ($1.47$, $37.5$ and $216$ at $w=0$, $0.1$, $1$), while
  \texttt{bsgs.negation} is $0.737$ $[0.609,0.892]$, below rho at every
  weight.
\item On two plane quartics, \texttt{bsgs.textbook} is $0.497$
  $[0.378,0.654]$ of plain rho.
\item \texttt{transfer.ghs\_cover} on two sextic-extension sources is
  $6.62$ $[5.13,8.54]$ of rho on the source, above at every weight.
\end{itemize}
The source classifies the study as a new capability, not an improvement:
at these sizes BSGS beats rho on the Jacobian, which is a statement about
small groups, and index calculus does not.

\paragraph{The Kummer collector.} A fixed-difference collector on genus-2
Kummer surfaces (the Scholten cohort: 21 models, 1{,}260 targets, 180 each at
16 to 40 bits) solved and verified every target in both its uncached and
cached arms, processing $2.30\cdot10^8$ quotient states
(\art{research/kummer_differential_collector_20261022/REPORT.md}). Both arms
lose in all 36 whole-cohort price cells and in all 45{,}360 target-cell
comparisons. Full-cold candidate/rho spans $8.96$--$55.94$ (uncached) and
$8.62$--$54.09$ (cached); at $S/M=0.5$, $I/M=10$ and $0.1$ GAE per native
operation the ratios are $14.36$ and $12.95$. The paired replay has a
passing validation gate (\art{benchmarks/20261022-kummer-paired-replay/}),
so these are validated measurements of a loss. On the way the study found
that the published Eq.~(11) of its Kummer recovery formula produced 1{,}257
incorrect signed divisors and three $D=0$ cases on the cohort; the corrected
generic equations recover them
(\art{research/kummer_recover_20261022/REPORT.md}).

\paragraph{Point-driven pencils.} The pencil collector builds Jacobian
relations from three-point pencils on the same frozen cohort
(\art{research/pencil_sparse_20261022/REPORT.md}). With three sparse
variable orders and a degree-peeling elimination, at $p=739$, $1{,}151$ and
$1{,}187$ (540 arm-target pairs over 180 targets), the native validator
recomputed all 19{,}440 target-level price cells and found 0 below rho. A
Markowitz pivot priority on the 60 targets at $p=739$ cut modular row updates
from 1{,}336{,}284 to 1{,}307{,}655 but raised logical index operations from
9.48 to 17.51 million, so its cold cost over peeling is $1.204855$ and 0 of 36
price cells are below either peeling or rho
(\art{research/pencil_markowitz_20261022/REPORT.md}). Both have passing
validation files (\art{benchmarks/20261022-pencil-degree-peeling/},
\art{research/pencil_markowitz_20261022/VALIDATION.json}). Later pencil
studies (a $p=739$ large-prime gate, early rejection, lifecycle and
inventory work) are engineering and report no new ratio.

\begin{table}[t]
\centering\small
\begin{tabularx}{\linewidth}{@{}>{\raggedright}p{3.6cm}>{\raggedright}X>{\raggedright\arraybackslash}p{3.1cm}@{}}
\toprule
Collector & Scope & Against matched rho \\
\midrule
genus-2 \texttt{ic.hyperelliptic} & two curves, 17.5 and 20.2 bits & $1.473\times$; $216\times$ at $w=1$ \\
Kummer fixed-difference collector & 21 models, 1{,}260 targets, 16--40 bits & $8.6$--$55.9\times$; 0 of 45{,}360 wins (validated) \\
point-driven pencils, peeling & $p=739$, 1{,}151, 1{,}187; 180 targets & 0 of 19{,}440 cells below (validated) \\
pencils, Markowitz priority & $p=739$, 60 targets & $1.205\times$ peeling; 0 of 36 below (validated) \\
genus-2 route at $n=2$ & 1{,}260 runs to 40 bits & $0.51$--$2.83$; unresolved \\
GHS cover from $\FF_{p^6}$ (ecbench) & two sources, 20.2 and 23.5 bits & $6.62\times$ \\
\bottomrule
\end{tabularx}
\caption[]{Measured Jacobian and cover collectors on prime-field or
small-degree sources. The $n=2$ genus-2 row is from
\S\ref{sec:ghs-odd}. Sources: as cited in the text.}
\label{tab:jac-samefield}
\end{table}

\paragraph{Other prime-field covers.} Three further results bound the
remaining constructions.
\begin{itemize}
\item \emph{Enge-type covers.} $C_{d,s}:y^2=f(t^d+s)$ covers $E$ with genus
  $(3d-1)/2$ for odd $d$; Enge's bound with $d=\Theta(\log p)$ gives an
  exponent of order $\log p\sqrt{\log\log p}$, larger than rho's
  $\frac12\log p$ (\art{research/enge_primefield_cover_20261024/REPORT.md}).
\item \emph{An external genus-2 backend.} Lindner's divisor-arithmetic
  library, pinned for comparison with the repository's Cantor arithmetic,
  returns the identity for \texttt{add(D,D)} on a valid divisor that is not
  2-torsion, while its \texttt{double(D)} returns a degree-2 divisor; its
  \texttt{add} cannot be used as a complete group law without an
  equal-input dispatch (\art{research/genus2_external_gate_20261021/README.md}).
\item \emph{Genus 3 over $\FF_{2^{83}}$.} A faithful transfer of the
  $n=83$ source into its genus-3 Jacobian image was replayed at toy sizes
  $n=9,13,19,23$ (400 of 400 destination attempts verified); only generic
  square-root solvers apply to an image of order $r\approx q$, and the
  ambient genus-3 index calculus loses to source rho by about 70.17
  exponent bits there (\art{research/n83_genus3_image_solver_20261031/README.md}).
\end{itemize}

\section{Joux--Vitse in November}\label{sec:jac-jv}

The Joux--Vitse cover~\cite{jouxvitse2012} of $E(\FF_{p^6})$, for curves with
full rational 2-torsion and $\#E=4\ell$, maps the problem to a genus-3
Jacobian over $\FF_{p^2}$. Chapter~\ref{sec:transfer} summarises the two arms
of this cover; this section is the full account: the validated sieve row at
$p=1511$, the November studies around it, and the earlier comparisons as
they now stand. All of them are on the weak class; none is on a prime-order
curve, and all but the validated row and the Hilbert stop below lack a
validation file.

\paragraph{The validated row: the sieve against a matched negation rho.} A
frozen protocol of 2 November re-measured the sieve arm
(\texttt{jv\_cover --exp sieve}) at $p=1511$ ($\ell=2^{61.4}$) against a
matched negation-map rho: 16 new rho runs pooled with 126 earlier ones under
a frozen homogeneity rule (142 runs), with the sieve rows of the 19 October
study unchanged and reproduced by a second build
(\art{research/jv_sieve_negrho_20261102/}). The sieve route costs $0.00647$
$[0.00545,0.00799]$ of rho in group additions (U1) and $0.0156$
$[0.0131,0.0192]$ in textbook $\FF_p$ multiplications (U2; seeds 17--24,
95\% bootstrap intervals). Rho is therefore $155\times$ $[125,184]$ and
$64\times$ $[52,76]$ more expensive. Charging the cover's unemitted F4,
descent and linear-algebra work at its upper bound leaves the ratio at
$0.107$--$0.123$ for $w=0.1$ and $1$. The order holds across the band in
all four cells (U1 and U2, emitted and upper-bound layers), for both seed
sets and both rho references. An independent validator passed the
comparison (\texttt{VAL-20261103-c31c76}), with these limitations:
\begin{itemize}
\item the replay is sampled: 1 of the 16 new rho runs, 1 of the 126 pooled
  runs and 2 of the 16 cover rows;
\item the cover's F4, descent and linear-algebra additions are not emitted,
  so only an upper-bound proxy prices them;
\item memory is not reported for either arm;
\item the run was local, not on the isolated lab.
\end{itemize}
The validator classifies the result as accounting, not an improvement to
any method: it is a matched measurement of an existing solver. It is the
programme's one validated measurement of index calculus below a matched rho.
It has no \texttt{pcc} claim record, so under the proof-carrying-claims rule
it is a validation-gate result whose statement is an unverified assertion.
It holds on the weak class only, says nothing about prime-order curves, and
is not extrapolated. A red-team review of the same study
(\art{research/jv_sieve_fbline_20261102/review/red_team/REPORT.md},
\texttt{RT-20261103-fd05ef}) adds two qualifications: U1 prices rho at a
schoolbook harness about $2.4\times$ above the tuned price, so the
like-for-like cell is U2 ($64\times$); and a minimal $\FF_{q^3}$ product would
make rho about $1.34\times$ cheaper in $\FF_p$ multiplications. Neither
reverses the order.

\paragraph{The Nagao arm near parity.} The 1 November crossover study ran
the plain arm at $p=1301$, 1451 and 1601 against two rho references on the
same instances (\art{research/jv_cover_crossover_20261101/README.md};
Table~\ref{tab:jac-nagao}):
against the distinguished-point rho it reads $0.793$ $[0.587,1.119]$,
$0.775$ $[0.645,0.970]$ and $0.587$ $[0.448,0.811]$, below at the two larger
primes under the frozen rule; against the strong rho, $1.349$
$[1.085,1.699]$ (above), $1.063$ and $0.910$ (not resolved). The source calls
this a candidate needing validation. A frozen re-scoring against a pooled
matched \emph{negation} rho (126 runs) reversed it
(\art{research/jv_cover_negrho_rescore_20261101/README.md}): in group
additions the cover is above rho at 1301 ($1.405$ $[1.206,1.629]$) and 1451
($1.160$ $[1.020,1.338]$) and unresolved at 1601 ($0.870$ $[0.697,1.183]$), and
in textbook $\FF_p$ multiplications it is above at every prime ($3.38$, $2.79$,
$2.09$).

\begin{table}[t]
\centering\small
\begin{tabular}{@{}llrrr@{}}
\toprule
Study & Reference, unit & $p=1301$ & $p=1451$ & $p=1601$ \\
\midrule
crossover & distinguished-point rho & $0.793$ & $0.775^{\downarrow}$ & $0.587^{\downarrow}$ \\
 & strong rho & $1.349^{\uparrow}$ & $1.063$ & $0.910$ \\
\addlinespace
re-scoring & negation rho, additions (U1) & $1.405^{\uparrow}$ & $1.160^{\uparrow}$ & $0.870$ \\
 & negation rho, mults (U2) & $3.376^{\uparrow}$ & $2.787^{\uparrow}$ & $2.090^{\uparrow}$ \\
\addlinespace
line search & negation rho, additions (U1) & $0.791^{\downarrow}$ & $0.659^{\downarrow}$ & $0.548^{\downarrow}$ \\
 & negation rho, mults (U2) & $1.900^{\uparrow}$ & $1.585^{\uparrow}$ & $1.318^{\uparrow}$ \\
\bottomrule
\end{tabular}
\caption[]{The plain Nagao arm of the Joux--Vitse cover against rho on the
same instances, cover/rho. $\uparrow$: interval above 1 (rho cheaper);
$\downarrow$: interval below 1; unmarked: not resolved. The re-scoring and
line-search rows share one pooled reference of 126 negation-rho runs. No
validation file. Sources: \art{research/jv_cover_crossover_20261101/},
\art{research/jv_cover_negrho_rescore_20261101/},
\art{research/jv_cover_fbline_20261102/}.}
\label{tab:jac-nagao}
\end{table}

\paragraph{A factor-base line search.} Choosing the factor-base line made
the plain arm $1.71\times$ $[1.54,1.89]$ cheaper ($S_{\rm line}/S_{\rm plain}=0.587$
$[0.530,0.651]$, as the model predicted, with every instance cheaper)
(\art{research/jv_cover_fbline_20261102/README.md}). Against the same
negation-rho reference the cover then reads $0.791$, $0.659$ and $0.548$ in
group additions, below rho at all three primes under all three estimators,
but $1.900$, $1.585$ and $1.318$ in textbook multiplications. In
multiplications the order does not hold across the band, and the cover's
native count omits its $\FF_p$ additions, so a ``below'' could narrow under
complete counts. The modelled crossovers move to $p^*\approx1185$ in
additions and $1799$ in multiplications, inside the range a 64-bit build can
test. There is no validation file, so this is a candidate improvement of
the index-calculus baseline, not a ranked one.

\paragraph{The line search on the sieve route.} A frozen controlled study of
2 November wired the same factor-base line search into the sieve route and
compared three arms with the plain route P on 48 paired instances at
$p\in\{1009,1511,1823\}$, the search's own cost charged
(\art{research/jv_sieve_fbline_20261102/README.md}). Pooled $S$ ratios
against P are $0.623$ $[0.562,0.689]$ for the Horner search H, $0.560$
$[0.500,0.631]$ for the exhaustive table T and $0.600$ $[0.535,0.674]$ for the
table at $n=32768$ (U); with additions and lookups charged, H and U stay
below P ($0.550$, $0.525$). An independent validator passed the paired
comparison (\texttt{VAL-20261103-594245}), but the validation gate is not
admissible, because the protocol names no rho reference, and there is no
\texttt{VALIDATION.json} or \texttt{pcc} claim. The red team
(\texttt{RT-20261103-fd05ef}) broke the claim for T: an uncounted integer
multiply per table lookup, once charged, puts T's band at $w=0.1$ at $1.164$,
worse than P; H and U are unresolved against each other. This study is the
cover-only half of the H and U arms below.

\paragraph{The sieve route on the same instances.} The 3 November study ran
three sieve arms (plain P, Horner line search H, table line search U) and a
matched negation rho on the same 24 instances at $p\in\{1009,1511,1823\}$,
removing the product of ratios that earlier comparisons used
(\art{research/jv_sieve_line_vs_rho_20261103/README.md};
Table~\ref{tab:jac-sieve}). Every cell is below rho, and the order holds
across the band in both cost layers, including the upper-bound layer that
prices the unemitted F4, descent and linear-algebra work as native
operations, where the ratios are still about $10\times$ below rho at $w=1$.

\begin{table}[t]
\centering\small
\begin{tabular}{@{}lrrrr@{}}
\toprule
Arm & $S$/rho (U2) & $S^+$/rho$^+$ (U2) & $S$/rho (U1) & Distance \\
\midrule
P (plain) & $0.0244$ $[0.0206,0.0292]$ & $0.0624$ & $0.0101$ & $\times41$ \\
H (Horner line) & $0.0154$ $[0.0135,0.0180]$ & $0.0344$ & $0.0064$ & $\times65$ \\
U (table line) & $0.0141$ $[0.0122,0.0167]$ & $0.0324$ & $0.0059$ & $\times71$ \\
\bottomrule
\end{tabular}
\caption[]{The Joux--Vitse sieve route against a matched negation rho on
the same instances, pooled over 24 instances at $p=1009$, 1511 and 1823, with
the per-prime reference the frozen homogeneity rule selected. U2: tuned rho
in $\FF_p$ multiplications; U1: group additions; $S^+$: with additions and
lookups charged on both sides. Distance: how many times cheaper than rho,
in U2. Source: \art{research/jv_sieve_line_vs_rho_20261103/README.md}.}
\label{tab:jac-sieve}
\end{table}

The ratio falls with $p$ in every arm (U, $S$/rho U2: $0.031$, $0.0115$,
$0.0081$), because rho grows as $\sqrt\ell\sim p^3$ faster than the route.
The source lists its limits: one rho run at $p=1823$ was killed by the time
rule and excluded, which biases that prime against the cover; the frozen
bootstrap understates rho-run uncertainty (a run-level bootstrap widens the
pooled U interval to $[0.0114,0.0182]$, no verdict changes); H against U is
not resolved per prime; the run was local; and there is no validation file.
It is accounting and distance, not an improvement.

\paragraph{Cutting the sieve's additions.} Two changes that leave every
line, hit, relation and logarithm bit-identical remove almost all of the
line enumeration's additions and 17--19\% of the sieve's. Over 48 paired
instances the measured $S^+$ ratio, new over old, is $0.621$ $[0.614,0.629]$
while $S$ moves by $0.990$; in the \texttt{bitop/v1} model the effect is
$0.880$ or $0.953$, depending on how linear algebra is priced, because $S^+$
prices an addition as a full multiplication
(\art{research/jv_sieve_adds_20261103/README.md}). The phase split at
$p=1511$ (line enumeration 25\% of multiplications, sieve 24\%, descent 42\%,
linear algebra 9\%; $S=0.0060$, $S^+=0.0219$) is in
\art{research/jv_cover_next_leads_20261102/LEADS.md}, which also names the
sieve, not the Nagao arm, as the cheapest route for the class.

\paragraph{Earlier comparisons, as they now stand.}
\begin{itemize}
\item The plain Nagao arm with a sparse-FGLM linear stage read
  $1.59$--$1.98\times$ a pooled rho constant at $p=1009$ on 24 October, closing
  as about $p^{-2}$ (\art{research/jv_cover_p_scaling_20261024/};
  \S\ref{sec:ladder3}). That was a lower bound against an unmatched
  constant; the matched negation rho above supersedes it, and the crossing
  it extrapolated near $p\approx1330$ does not hold against that rho.
\item The 19 October sieve study (\art{research/jv_sieve_strong_rho_20261019/})
  measured $0.024$ $[0.019,0.031]$ of a strong rho at $p=1009$ and $0.0042$ at
  $p=1823$, but its 20 October review withheld every comparison (no numeric
  acceptance for the controls, uncharged rho hashes, incomplete accounting).
  Its sieve rows are the ones the validated row re-scored.
\item The replication of the 19 October sieve study failed its preregistered
  acceptance rule A4 (a rho A/A ratio within $[0.85,1.15]$: $0.792$
  $[0.572,1.073]$), so all its decisions are withheld; its descriptive
  sieve/rho ratios ran from $0.647$ at $p=503$ to $0.00404$ at $p=1823$
  (\art{research/jv_sieve_strong_rho_v2_20261020/README.md}). The registered
  prospective re-run (\art{research/jv_sieve_prospective_20261023/}) never ran
  and was superseded by the matched study above.
\item Re-priced in \texttt{bitop/v1} with both arms executed, the whole-route
  row (an isogeny walk from a non-weak challenge curve to the weak class,
  then the cover) at $p=1511$ ($\ell=2^{61.4}$) reads $0.0098$
  $[0.0069,0.0144]$ of the strong rho R3 (largest band upper bound $0.075$);
  it has no validator report
  (\art{research/ic_acct1_cover_bitop_20261021/README.md}).
\item A 9 October whole-route replay from a non-weak challenge curve, through
  an isogeny walk to a weak curve and its cover, read $0.058$ of a pooled,
  unmatched rho at $p=1009$ on one seed per prime, with an earlier revision of
  the solver whose source includes the sieve module; which arm ran is not
  recorded (\art{research/jv_cover_e2e_30bit_20261009/RESULT.md}). The
  cover-decomposition note still states a measured sieve crossover near
  $p\approx430$
  (\art{research/notes/index-calculus/RESEARCH_COVER_DECOMPOSITION_LEDGER.md}
  \S11) with no erratum; it should be read through the review. Leaderboard
  Table~D, the historical study, is withheld.
\item On the method side, the sparse-FGLM stage made the plain arm
  $1.28\times$ cheaper end to end
  (\art{research/jv_cover_sparse_fglm_20261018/}), and relation collection is
  $99.99\%$ of $S$ (\art{research/ic_cover_leads_20261018/}).
\item A large-prime variation of the cover decomposition does not pay at
  $p\le12{,}000$: linear algebra is $2.5\cdot10^{-3}$ of $S$ there, and one large
  prime saves at most $0.02\%$ (\art{research/jv_cover_large_prime_20261018/README.md}).
\end{itemize}

\paragraph{The Hilbert early stop, and a relabelling.} In the plain arm's
F4 stage, a Hilbert-function stop certifies that the six quadrics' leading
forms are a regular sequence and drops every pending pair that provably
reduces to zero (\art{research/ic_cover_weakclass_20261023/REPORT.md}).
Plumbed into the logarithm mode and measured with a new native counter, it
cuts counted $S$ by $1.848$--$1.875\times$ over 25 cells at equal relations,
with the emitted native counter at $0.983$--$0.990$ of the baseline:
engineering, with a passing validation gate
(\art{benchmarks/20261029-jv-cover-hilbert-stop/}). The form first found,
B0, cut counted $S$ $2.53$--$2.64\times$ while the unpriced F4-driver work rose
$8.3$--$8.5\times$; the study classifies B0 as relabelling, which is in the
corrections chapter (Table~\ref{tab:corr-transfer}). The baseline's distance
to rho falls with $p$: $S$/rho of $1361$, $205$, $35.2$ and $7.19$ at $p=53$,
101, 251 and 503, slope $-2.27\pm0.13$ in $\log p$ over four sizes, against
the matched negation rho of the Hilbert-stop study; the pooled-constant
series of Chapter~\ref{sec:frontier} reads $620$ at $p=53$ and slope $-2.05$.

\paragraph{The weak class.} Every number above is on curves with full
rational 2-torsion and $\#E=4\ell$. A norm-one conductor theorem
characterises the weak isogeny classes; a random full-2-torsion curve lies
in one with probability $0.5908$ $[0.5754,0.6059]$ at $p=11$ and $0.6320$
$[0.6169,0.6468]$ at $p=47$, over a census to $p\le53$
(\art{research/iso1_weak_classes_20261007/REPORT.md}). Generalising the cover
over 2-torsion configurations gives genera $g\in\{1,3,7,9,13\}$, and
Joux--Vitse's is complete among such covers with $g\le3$
(\art{research/jv_weak_generalization_20261008/README.md}). Depth-1 isogeny
classes hold no weak curve over 4{,}871{,}062 curves, while 86--95\% of deeper
classes hold one (\art{research/jv_weak_class_law_20261007/}). Since weakness is not a
class invariant, a Teske-type trapdoor in this class is structurally
possible; the atlas lead is unexamined.

\section{What the Jacobian routes settle}\label{sec:jac-summary}

\begin{itemize}
\item On prime fields every same-field Jacobian route is closed by
  theorem, and every measured collector (Kummer, pencils, genus-2 index
  calculus) is above rho, three of them under a passing validation gate.
\item At $n=3$ the quartic routes are below rho in the weak family under
  counted costs; on generic quartics the Diem arm is below at $w=0$ only.
\item At $n=4$ and $5$ the Jacobians exist in the right classes; the one
  complete attack (genus 5) loses at every $q$ measured.
\item The Joux--Vitse sieve is between about 16 and 71 times cheaper than a
  matched negation rho on the weak class at $p\le1823$ in $\FF_p$
  multiplications, depending on arm and on whether additions are charged,
  and $155\times$ in group additions ($64\times$ in $\FF_p$ multiplications)
  at $p=1511$; only the $p=1511$ comparison is validated, and it has no
  \texttt{pcc} claim. The plain Nagao arm is above a matched negation rho at
  $p=1301$ and 1451 and unresolved at 1601 in group additions, and above at
  all three in multiplications; with the line search it is below in
  additions only, unvalidated.
\end{itemize}
\chapter{Extension fields: the GLV-invariant programme and quotients}\label{ch:glvprogramme}

Chapter~\ref{ch:extension} gives the $D_4$ system at $k=4$, its engineering
ladder, the characteristic-two distance to rho and the red-team verdict on
its crossover; \S\ref{sec:gaudry} gives the first $D_4$ build and the
modelled handover it moved. A second line of work ran beside them: one
note, \art{research/notes/index-calculus/RESEARCH_GLV_INVARIANT_FACTOR_BASES.md}
(``the GLV note'' below; \S\S1--9.57, about 6{,}000 lines), registered and
measured experiments E1 to E42 on factor bases invariant under
automorphisms, Frobenius-type maps and 2-torsion translations, first on
subfield and GLS curves, then at $k=4$ over $\FF_{p^4}$. The extension
chapter quotes one line of it. This chapter gives the programme's results
in order, the corrections it made to its own reference, and the quotient
and cofactor levers that belong with it. Validation is named per
experiment; the rest is producer-reported.

\section{Folds on subfield and GLS curves: E1--E18}\label{sec:glv-folds}

\paragraph{What a fold can be.} E1--E7 established one closed form
(GLV note \S6.8). The fold of a factor base by any finite group of
endomorphisms is the order of the subgroup of $(\ZZ/r\ZZ)^*$ that the group's
eigenvalues generate: 12 on the $j=0$ GLS twist, 4 (not 8) on the $j=1728$
twist, 3 (not 9) on the $j=0$ subfield curve. Rho folded by the same group
takes the square root of that number fewer steps, verified (E6, E9: steps
$\div1.40$, $1.74$ and $2.16$ at $A=4$, 6 and 12 against the expected $1.41$,
$1.73$ and $2.45$). CM endomorphisms of degree 2 and 3 and Q-curve
endomorphisms keep base points in the base only at chance (E4, E14), and
where a cofactor is fixed by the maps, two-summand relations stay inside one
orbit, so the fold buys columns and not relations (E5). Every class in
E1--E16 is engineering or accounting.

\paragraph{The best subfield arm and its boundary.} With three summands on
the Frobenius line, FGHR's 2-torsion symmetry on the system (a degree-16
resultant in place of a 64-dimensional quotient, $S\div27.0$) and the
translation fold on the base ($S\div2.52$ more), the cheapest arm was
$267$--$5{,}741$ times rho (E13). E16 then stated why no lever on these
families can do better than a constant: on a subfield curve the large prime
lies in the trace-zero part, $r\approx p^2$, so rho costs $\Theta(p)$ and a
line base needs $\Theta(p)$ relations; to $p=2^{16}$ the best arm's ratio to
rho grows as $r^{0.18}$ ($267$--$6{,}616\times$). E16 also corrected two
earlier exponent readings, E11's and E13's.

\paragraph{Moving the levers to $k=4$.} E17 carried the $D_3$ solve
($\div38.0$) and the $\tau_T$ fold ($\div2.00$ exactly) to the full group
$E(\FF_{p^3})$, still $180$--$377\times$ rho. E18 folded the $k=4$ linear
algebra: on $j=0$ curves the $\psi$ fold takes the asymptotic ratio to the
matched walk from $1.232$ to $0.138\pm0.008$ and the extrapolated handover
from $2^{156.1}$ to $2^{146.0}$, the figure \S\ref{sec:gaudry} quotes; the
$\tau_T$ fold cannot pay for its curves' cofactor $h=4$ and $4.7\times$
poorer yield. E18 also re-read an earlier note's ratio $0.518$ as $0.733$
against the negation walk (GLV note \S9.4).

\section{The $V_4$ line and the $k=4$ baseline: E19--E26}\label{sec:glv-v4}

\paragraph{E19: the first end-to-end run.} The $\tau_T$/$D_4$ pipeline
recovered logarithms at $k=4$ with every phase priced, on two curves each at
$p=53$--$269$ ($n=2^{20.9}$--$2^{30.3}$, $h=4$); the solves are $100.0\%$ of
$S$ at a median $1.37$--$1.41\cdot10^9$ $\FF_p$ multiplications per call, and
the method is $2.6$--$8.8\cdot10^6$ times the negation rho (GLV note \S9.6).
From the pipeline's own constants the modelled handover is $2^{116.7}$,
reproducing the $2^{119}$ assembled from separate harnesses. E20 found that
double large primes have the sub-rho exponent at $k=4$ and lose by the
constant ($1.66$--$3.60\times$ the full base, crossing it only near
$2^{164}$--$2^{222}$), and E21 that trace replay on the weighted solve is a net
loss ($0.95$--$0.98\times$ with fallbacks charged), because the rows that
vanish are 7\% of the matrix and differ from target to target.

\paragraph{E22: the $V_4$ line.} On curves over $\FF_{p^4}$ with full
rational 2-torsion and rational halves of two 2-torsion points, a line
exists that is stable under the whole Klein group $V_4$ of translations, and
$\Sem5$ can be written in $V_4^4\rtimes S_4$ invariants: the quotient falls
from 512 to 64 and the Macaulay degree from 26 to 11. One solve costs about
$1.7\cdot10^6$ multiplications, $\approx820\times$ less than the weighted $D_4$
solve, but the larger translation orbits and the thinner base cost about
$10\times$ in targets per relation (GLV note \S9.12). Its failures are all
planted targets with two points of one translation orbit, which give a
positive-dimensional system; uniform targets meet this with probability
$\approx\mathrm{cols}^2/n$.

\paragraph{E24: the baseline moves, validated.} End to end on twelve curves
at $p=101$--$521$ ($n\le2^{34.1}$), the $V_4$ pipeline recovered every
logarithm; the solves are $99.3$--$99.5\%$ of $S$ and every row is below E19
at matched $p$ by $35$--$240\times$ (GLV note \S9.14; Table~\ref{tab:glv-ladder}).
Table~\ref{tab:glv-e24} gives the rows. The fitted slope of $S$/rho in $n$ is $-0.25$, the theory's $n^{-1/4}$, and the
modelled handover moves from E19's $2^{116.7}$ to $2^{90.1}$ ($2^{91.7}$ from
the largest $p$ alone). The class is engineering, a constant, and under the
success criterion a verified, measured improvement of the baseline. An
independent validator replayed every run of E19, E24 and E25 field for field
and recomputed the claims (\art{benchmarks/20261018-k4-v4-line/},
\art{benchmarks/20261018-k4-gls-v4/}; both gates pass). It found seven
textual errors, none changing a result (\S9.19): for example the
E19-to-E24 factor at $p=101$ seed 2 is $170\times$, not $171\times$; 392 kept
rows were re-added, not 381; and unsolved targets reach $0.033\%$ per curve,
not $0.012\%$. It also priced the band: the $35$--$240\times$ is at $w=0$; with
one GAE (97 multiplications) per Macaulay entry and walk step it is
$19$--$132\times$ at $w=0.1$ and $6.5$--$45\times$ at $w=1$, and E24 stays below
E19 on every curve at every weight. Every E24 row is still
$1.3\cdot10^4$--$1.7\cdot10^5$ times rho, and rho is ahead across the band.

\begin{table}[t]
\centering\small
\begin{tabular}{@{}rrrrrrr@{}}
\toprule
$p$ & seed & $\log_2n$ & base/cols & $C$ median & $S$/rho & E19/E24 \\
\midrule
101 & 1 & 24.6 & 36/9 & $1.62\cdot10^6$ & $1.69\cdot10^5$ & $37\times$ \\
101 & 2 & 24.6 & 52/13 & $1.62\cdot10^6$ & $3.43\cdot10^4$ & $170\times$ \\
149 & 1 & 26.9 & 64/16 & $1.63\cdot10^6$ & $3.80\cdot10^4$ & $108\times$ \\
149 & 2 & 26.9 & 60/15 & $1.63\cdot10^6$ & $5.54\cdot10^4$ & $61\times$ \\
197 & 1 & 28.5 & 76/19 & $1.65\cdot10^6$ & $1.40\cdot10^5$ & $35\times$ \\
197 & 2 & 28.5 & 84/21 & $1.64\cdot10^6$ & $5.44\cdot10^4$ & $52\times$ \\
269 & 1 & 30.3 & 132/33 & $1.65\cdot10^6$ & $2.30\cdot10^4$ & $240\times$ \\
269 & 2 & 30.3 & 148/37 & $1.65\cdot10^6$ & $2.34\cdot10^4$ & $110\times$ \\
397 & 1 & 32.5 & 208/52 & $1.66\cdot10^6$ & $1.31\cdot10^4$ & \\
397 & 2 & 32.5 & 188/47 & $1.66\cdot10^6$ & $1.58\cdot10^4$ & \\
521 & 1 & 34.1 & 240/60 & $1.66\cdot10^6$ & $1.88\cdot10^4$ & \\
521 & 2 & 34.1 & 232/58 & $1.66\cdot10^6$ & $2.15\cdot10^4$ & \\
\bottomrule
\end{tabular}
\caption[]{E24: the $k=4$ pipeline on the $V_4$ line, $h=4$ on every curve,
against the matched negation rho of the time ($w=0$). $C$: $\FF_p$
multiplications per solve call. E19/E24: the earlier pipeline's $S$/rho over
this one at the same $p$ and seed, on different curves of that $p$; E19 did
not run $p>269$. The $p=101$ seed~2 factor is the validator's $170\times$
(printed $171\times$). Validated: \art{benchmarks/20261018-k4-v4-line/}.
Source: GLV note \S9.14 and \S9.19.}
\label{tab:glv-e24}
\end{table}

\paragraph{E25: the GLS fold.} On GLS twists over $\FF_{p^4}$ the GLS map
$\psi$ ($\psi^2=-1$) keeps the $V_4$ line and acts on it as $z\mapsto\pm i/z$, so
the base folds by $\langle V_4,\psi\rangle$ with 8 points a column (GLV note
\S9.16). Measured on the same curve, base and target stream, the fold cuts
$S$ by $1.35$--$2.67\times$ (median $1.86$), and against the matched
$\langle\pm\psi\rangle$ walk the $\psi$ arm gains $1.26$--$1.97\times$ over the
control at every $p\ge197$; below that, walk set-up dominates. The modelled
handover on these curves moves from $2^{93.2}$ to $2^{88.0}$
(Table~\ref{tab:glv-e25}).

\begin{table}[t]
\centering\small
\begin{tabular}{@{}rrrrrrr@{}}
\toprule
$p$ & seed & $\log_2n$ & cols A/B & $S_A/S_B$ & $S_B$/rho$_{\pm\psi}$ & improvement \\
\midrule
101 & 1 & 22.6 & 10/5 & 1.94 & $6.77\cdot10^4$ & 2.26 \\
101 & 2 & 22.6 & 14/7 & 2.11 & $2.93\cdot10^4$ & 2.05 \\
149 & 1 & 24.9 & 10/5 & 2.67 & $2.65\cdot10^5$ & 3.13 \\
149 & 2 & 24.9 & 14/7 & 1.77 & $7.01\cdot10^4$ & 2.25 \\
197 & 1 & 26.5 & 22/11 & 1.35 & $6.11\cdot10^4$ & 1.26 \\
197 & 2 & 26.5 & 30/15 & 1.69 & $2.62\cdot10^4$ & 1.51 \\
269 & 1 & 28.3 & 30/15 & 1.74 & $4.84\cdot10^4$ & 1.54 \\
269 & 2 & 28.3 & 40/20 & 2.02 & $1.56\cdot10^4$ & 1.44 \\
397 & 1 & 30.5 & 60/30 & 1.72 & $1.13\cdot10^4$ & 1.76 \\
397 & 2 & 30.5 & 48/24 & 2.32 & $1.71\cdot10^4$ & 1.97 \\
\bottomrule
\end{tabular}
\caption[]{E25: the GLS fold on the $V_4$ line, $h=16$ on every curve. Arm
A: $V_4$ fold, against the negation walk; arm B: $\langle V_4,\psi\rangle$ fold,
against the $\langle\pm\psi\rangle$ walk, on the same curve, base and target
stream. Improvement: $(S_A/\text{rho}_{\rm neg})/(S_B/\text{rho}_{\pm\psi})$.
Distances predate E27's stronger walk. Validated:
\art{benchmarks/20261018-k4-gls-v4/}. Source: GLV note \S9.16.}
\label{tab:glv-e25}
\end{table} It does not move
E24's rows: every curve of this construction has $h=16$, which makes rho
$2\times$ cheaper at a given $p$ and absorbs most of the fold at matched $n$.
E26 then showed, on 560 GLS twists, that $h=16$ is forced: the halves
condition for the $V_4$ line holds exactly on the 250 twists with $4\mid t$,
each with $v_2(\#E)\ge4$, and fails on all 310 with $t\equiv2\pmod4$ (\S9.18).
The validator noted that the $\langle\pm\psi\rangle$ walk applies $\psi$ to both
coordinates on every step and estimated that comparing abscissae first would
save about $6\%$, so E25's distances lean that much toward index calculus.

\paragraph{E23 and E28: the other lines.} On $j=0$ the order-24 group exists on
the $x_0\FF_p$ line, but the line carries no point of the large subgroup (E23,
by derivation, confirmed on points). On E19's single-$\tau_T$ line with GLS
twists, $\psi$ keeps the line with $h=4$ on all 40 curves and halves the
columns, but the solve is unchanged at $\approx1.4\cdot10^9$, so the line stays at
least $6.99\cdot10^5$ times rho in a model (E28, \S9.30). The validator refused E28 at
the gate because its end-to-end $S$ is a model.

\section{The solve and its pivots: E31--E42}\label{sec:glv-chain}

With the relation phase priced and the base folded, the per-call cost of the
degree-11 solve is all of $S$. Six experiments worked on it.

\paragraph{E31: the chain solve.} The quotient of the $V_4$ system is a
12-generator module over $\FF_p[P_2]$, but the solver paid for a dense 64-wide
normal form at each of 240 pivots, a dense Hessenberg reduction and a dense
64-dimensional kernel (GLV note \S9.23). Working in the module (left-looking
normal forms of the 12 chain ends, $\det R(x)$ of a $12\times12$ polynomial
matrix, kernels of $R(\lambda)$) cut one call from $1.62$--$1.66\cdot10^6$
multiplications to $3.04$--$3.33\cdot10^5$ with the sparsest-pivot rule (B),
with the same decompositions on all 10{,}000 calls; end to end on E25's curves
$S$ fell $4.53$--$4.76\times$ (B) and $2.50$--$2.56\times$ (A, without the pivot
rule), with every relation, logarithm and walk unchanged. The independent
validation \texttt{VAL-20261024-24ef8b} passed
(\art{benchmarks/20261024-e31-chain-solve/}) with seven errata and one
qualification that changed the ranking (\S9.24): B's sparsest-pivot rule
counts the nonzeros of every row it builds, which was neither counted nor
priced; priced at the band's rate, B falls behind E25 at $w=0.1$ and $1$, so
\textbf{B beats E25 and A only at $w=0$}, its rank is unresolved across the
band, and A, not B, is the ranked solver. A stays below E25 at every weight
($2.50$--$2.56\times$ at $w=0$, $1.39$--$1.43\times$ at $0.1$, $1.034$--$1.053\times$
at $1$): at $w=1$ the $393\times448$ matrix the chain solve still builds is almost
all of the cost.

\paragraph{E33--E39: pivot rules under the band.} The next experiments asked
which pivot rule pays for its own selection work.
\begin{itemize}
\item \emph{E33.} Counting B's selection pass gave $1.05\cdot10^6$ native
  operations per call, about six matrix reads; at $w=0.1$ B falls to $0.19\times$
  A. A counted estimate B$''$ with $O(1)$ work per row update is ahead of A at
  $w=0$ ($1.54$--$1.56\times$) and $0.1$ ($1.046$--$1.048\times$) and behind at $1$:
  not ranked (\S9.37).
\item \emph{E35.} Replaying B's pivot order per curve does not work: no two
  calls of a curve share a sequence, because ties are broken by accidental
  zeros of the target's coefficients; the replay keeps $18.8$--$36.8\%$ of B's
  saving and splits 14 to 6 at $w=1$: unranked (\S9.42). This agrees with E21's
  finding on a different object.
\item \emph{E36--E38.} Taking each bucket's first-arrived row (F) needs no
  selection work and cut A's per-call cost $1.47$--$1.55\times$; F and two
  variants were ahead of A and E25 at every weight on all 20 rows, then on all
  32 rows on a fresh stream. Each time a registered control's kill criterion
  fired (on an inverse-map control in E36, on an F-specific control's positive
  check in E37, at the pilot in E38, because every column the normal forms
  read is contested), so none was ranked (\S\S9.43--9.47).
\item \emph{E39.} On the third attempt, with controls whose positive checks
  see what each defect corrupts, no kill fired, and F, S and H are ranked
  ahead of A and E25 at $w=0$, $0.1$ and $1$ on all 32 rows: per row
  $1.46$--$1.54\times$ fewer solve multiplications for F at $w=0$ and
  $S_A/S_F=1.0115$--$1.0130$ at $w=1$. It is validated (31 of 31 runs replayed,
  1{,}037{,}565 fields identical, 7{,}286 summary values recomputed with 0
  disagreements; \texttt{VAL-20261102-E39},
  \art{benchmarks/20261102-e39-pivot-rules/}).
\end{itemize}
The sequence is an example of the validation gate doing its job: a measured
order held in three runs and was ranked only when a preregistered run with
admissible controls held.

\paragraph{E40 and E42b: the matrix itself.} At $w=1$ the band is set by the
Macaulay entries (176{,}064 per call), which no pivot rule removes. E40 found
that the forward echelon and the normal forms never read the 73{,}237 cells left
of each row's lead (41.6\%), and stored each row from its lead: $1.71\times$ fewer
cells, with A's arithmetic, pivots and outputs bit for bit (\S9.51). Ranked by
its registered rule, the stored form V equals A at $w=0$ and is ahead by
$1.25$--$1.26\times$ at $0.1$ and $1.36$--$1.38\times$ at $1$ on all 32 rows; it has a
passing gate (\texttt{VAL-20261102-E40},
\art{benchmarks/20261102-e40-macaulay-entries/}). It is the first change at the
$w=1$ end since E31, and an improvement of the baseline at the positive weights
only. E42 stopped on a mis-registered growth-control count; its re-registration
E42b reserved V's rows to the block after their build end, cutting cells per
call from 102{,}827 to about 92{,}970 at about 5{,}500 charged checks, equal to V at
$w=0$ and ahead by $1.022$--$1.023\times$ and $1.032$--$1.034\times$ at $0.1$ and $1$
(\S9.57); no validation is recorded for it.

\section{The reference, corrected: E27, E30 and E32}\label{sec:glv-reference}

\paragraph{E27: the matched walk was too expensive.} The rho reference of
E24--E31 used 19-multiplication $\FF_{p^4}$ products and one $\FF_p$ inversion per
addition. A Karatsuba tower and Montgomery-batched norm inversions cut the
affine addition from 97 to $66.3$--$66.7$ multiplications and made the same
walks, with the same points, collisions and logarithms, $1.4533$--$1.4634\times$
cheaper (GLV note \S9.29). The index-calculus side moved only
$1.0048$--$1.0057\times$, because its $\FF_{p^4}$ work is about 2\% of $S$. So every
row moved about $1.45\times$ \emph{further} from rho at $w=0$ ($1.24$--$1.33\times$ at
$w=1$), to $6.39\cdot10^3$--$4.06\cdot10^5$ times rho. The correction is validated
(\texttt{VAL-20261031-e27a1c}, \art{benchmarks/20261030-e27-walk-karatsuba/})
and is also recorded with the other reference corrections
(\S\ref{sec:rv-aut}, Table~\ref{tab:corr-rho}). Replaying E25's walks with the
current code also exposed that a later commit had begun charging the
$[d]G=Q$ check inside the walk ($0.13$--$0.78\%$ of a walk), raising every matched
rho recorded since; E27 prices walks without it, as E25 did.

\paragraph{E30: an accounting audit.} Trivial products are $0.31\%$ of $S$, so
the distance to rho moves by $-0.31\%$ at most; validated (GLV note \S\S9.32--9.33).

\paragraph{E32: the handover ladder, validated.} Instead of a single-point
model, E32 measured the distance over eleven sizes, $p=101$--$3089$, three
curves each ($n=2^{22.6}$--$2^{42.4}$, 33 paired curves, 528 walks), on the line's
best configuration (arm B with solver A and Karatsuba) against its strongest
matched walk (\S9.36). The per-size median distance falls from $2.1\cdot10^4$ to
$1.5\cdot10^3$ at $w=0$ (best row $1.14\cdot10^3$ at $p=3089$), and is $4.9\cdot10^3$ at
$w=0.1$ and $2.0\cdot10^4$ at $w=1$ at the largest size. The crossover is empty in
the reachable range: $n^*>2^{42.4}$, measured. The free-slope fit gives
$\beta=-0.243$ with an interval of $\pm0.04$ containing $-1/4$, and a fitted
crossover at $2^{85.5}$ $[2^{79.2},2^{94.3}]$ at $w=0$, $2^{91.2}$
$[2^{84.4},2^{100.7}]$ at $0.1$ and $2^{97.6}$ $[2^{90.1},2^{107.9}]$ at $1$. The
15-bit width of the interval is what an extrapolation over 43 bits beyond the
top size can say. The ladder has a passing gate
(\art{benchmarks/20261031-e32-handover-ladder/}); its validator corrected four
printed statements without changing a verdict.

\begin{table}[t]
\centering\small
\begin{tabularx}{\linewidth}{@{}l>{\raggedright}X>{\raggedright}p{3.3cm}>{\raggedright\arraybackslash}p{2.6cm}@{}}
\toprule
Step & Change & Effect & Gate \\
\midrule
E19 & $\tau_T$/$D_4$ pipeline end to end & $2.6$--$8.8\cdot10^6\times$ rho; handover $2^{116.7}$ (model) & validated \\
E24 & $V_4$ line, invariant solve & $35$--$240\times$ below E19 ($w=0$); handover $2^{90.1}$ (model) & validated \\
E25 & GLS fold on the $V_4$ line ($h=16$) & $S\div1.35$--$2.67$; handover $2^{88.0}$ (model) & validated \\
E31 & chain solve, solver A & $S\div2.50$--$2.56$ ($w=0$), $\div1.034$--$1.053$ ($w=1$) & validated; B unranked \\
E27 & reference: Karatsuba, batched inversions & every row $1.45\times$ further from rho; $6.39\cdot10^3$--$4.06\cdot10^5\times$ & validated \\
E32 & distance ladder to $n=2^{42.4}$ & $1.5\cdot10^3\times$ at the top; fitted $n^*=2^{85.5}$ $[2^{79.2},2^{94.3}]$ & validated \\
E39 & arrival-order pivots F & $1.46$--$1.54\times$ fewer solve mults ($w=0$); $1.01\times$ ($w=1$) & validated \\
E40 & rows stored from their lead & equal at $w=0$; $1.25\times$ ($0.1$), $1.36$--$1.38\times$ ($1$) & validated \\
\bottomrule
\end{tabularx}
\caption[]{The $k=4$ line over $\FF_{p^4}$ in the GLV note, as ranked. Effect:
against the previous incumbent unless stated; distances are to the matched
rho of the time, which E27 strengthened. Every row has rho ahead at every
weight of the band; every handover is a model or a fit. Sources: GLV note
\S\S9.6--9.57; validation directories under \art{benchmarks/}.}
\label{tab:glv-ladder}
\end{table}

Two things change the reading of the earlier single-point figures. The
modelled handovers for odd characteristic, $2^{135.1}$--$2^{135.8}$ in
\S\ref{sec:gaudry} and the $2^{119}$ of the GLV note's early frames (E19,
\texttt{d4-s5}), are
superseded on this family by a measured distance to $2^{42.4}$ and a fitted
interval of $2^{79}$--$2^{108}$ across the band. And the family is narrow: GLS
twists over $\FF_{p^4}$ with $p\equiv1\pmod4$, full rational 2-torsion, the halves
condition and $h=16$. A curve without rational 2-torsion, such as a
prime-order curve, has none of the symmetry this line uses.

\section{The order-24 twist action: E29, E34 and E41}\label{sec:glv-order24}

On E18's $j=0$ curves over $\FF_{p^4}$, the Frobenius-derived map
$\psi_p=(u^2x^p,u^3y^p)$ with $u=b^{-(p-1)/6}$ has eigenvalue of order 24 on all
eight curves ($p=241$--$1069$), and with the automorphisms it acts as a cyclic
group of order 24 on $\langle G\rangle$ (GLV note \S9.31). It changes two things.
\begin{itemize}
\item \emph{The base.} On the canonical model the order-24 fold divides the
  columns by $4.00$ and the linear algebra by $13.3$--$15.2$ against E18's cubic
  fold, and moves the modelled $\Sem5$ handover by about 9--10 bits, to
  $2^{135}$--$2^{142}$. No $\Sem5$ solve was run, so this is a model; the
  validator's report for E29 is a fail at the gate for that reason alone,
  with every measured part verified.
\item \emph{The walk.} An order-24 walk takes $1.70\times$ fewer steps than an
  order-6 walk (pooled), but E29's class representative cost $191$--$210$
  multiplications a step against $109$--$122$, so the walk's cost ratio was
  $0.962$ pooled at $w=0$. The validation added that with one native
  operation per step priced the ratio is $0.996$ at $w=0.1$ and $1.205$ at $w=1$,
  so neither walk was the stronger reference across the band.
\end{itemize}
E34 built a cheaper representative, 15.17 multiplications instead of 86--92,
so an order-24 step costs 7--10\% more than an order-6 step instead of 80\%
(\S9.39). Pooled over the curves, E29's walk over the cheap one is $1.563$,
$1.254$ and $1.126$ at $w=0$, $0.1$ and $1$ (schoolbook), order 24 ahead on 7 of 8
curves; the cheap order-24 walk became the matched rho on E18's $j=0$ curves,
and E18's $\psi$ rows moved $1.13$--$1.63\times$ further from rho
(\texttt{VAL-20261101-146a90}). E41 then asked why the order-24 walk fell short of
its ideal factor of 2 (\S9.53). The shortfall was sampling noise in E34's 16-seed
design: on fresh seeds, 1{,}024 walks per arm, the order-6 walk sits at $0.969\pm0.016$
and the order-24 walk at $1.006\pm0.016$ of their ideals, and no order-24-specific
cause exceeds $0.1\%$ of steps. The corrected reference is resolved at every
weight against both order-6 walks; E18's $\psi$ rows move $1.761$, $1.418$ and
$1.275\times$ from rho at $w=0$, $0.1$ and $1$ against E29's walk, and the GLS rows
of Table~\ref{tab:glv-ladder} do not move (\texttt{VAL-20261102-918856}). All three
are reference corrections, not index-calculus improvements; E29's base fold
leaves the $j=0$ line $\ge2^{45}$ behind the $V_4$ line's modelled handover.

\section{Quotients and cofactors}\label{sec:glv-quotients}

Quotienting the problem by a symmetry and clearing a cofactor are the
other two ways an extension-field structure enters a factor base.
\begin{itemize}
\item \emph{Symmetry quotients of the system.} The EcGFp5 quotient by
  $(\ZZ/2)^4\rtimes S_5$ at $m=5$ and its block-Wiedemann finish are in
  \S\ref{sec:pt-gfp5}; they are per-system baselines over toy primes and are
  not repeated here.
\item \emph{Shared systems for cofactor lifts.} On Curve25519-family
  miniatures at $k=20$, solving one shared F4 system for all eight cofactor
  lifts uses one call per target instead of eight but costs $30.7$--$628\times$ the
  multiply-adds of separate lifts, because the maximum solving degree rises
  from 5 to 8--10 (\art{research/curve25519_cofactor_shared_20261102/README.md});
  the Curve25519 programme is in \S\ref{sec:c25519}.
\item \emph{Cofactor-cleared bases.} On prime curves with $h=32$ the cleared
  base's A/B test returned \texttt{invalid\_controls} by its frozen rule and is
  not cited as a gain (\S\ref{sec:pg-cofactor}); the $m=2$ torsion-anchor
  reading at $n=37$ was reversed on its bootstrap recheck
  (\S\ref{sec:corr-binary}).
\item \emph{A relabelled quotient count.} In the Joux--Vitse cover, the
  first form of the Hilbert early stop (B0) cut counted $S$ $2.53$--$2.64\times$
  while unpriced F4-driver work rose $8.3$--$8.5\times$; the study classifies it as
  relabelling, and only its repair is engineering (\S\ref{sec:jac-jv},
  Table~\ref{tab:corr-transfer}).
\end{itemize}
The common lesson is the one E21 and E35 drew on the $V_4$ solve: a quotient
or a shared system pays only if the work it moves into an unpriced counter is
priced, and the band is where that shows.

\section{What the extension-field programme settles}\label{sec:glv-summary}

\begin{itemize}
\item A fold's size is the order of the eigenvalue subgroup; on subfield
  curves no fold can do better than a constant, because $r\approx p^2$.
\item At $k=4$ on GLS twists with full rational 2-torsion, the baseline moved
  by validated constants: the $V_4$ line ($35$--$240\times$), the GLS fold, the
  chain solve, arrival-order pivots and lead-stored rows. The reference moved
  the other way ($1.45\times$), and the measured distance to rho is about
  $1.5\cdot10^3$ at $n=2^{42.4}$, falling as $n^{-0.24}$.
\item The crossover is a fit, $2^{85.5}$ $[2^{79.2},2^{94.3}]$ at $w=0$ and later at
  higher weights, on a family no deployed curve belongs to.
\item On $j=0$ curves the order-24 action is the strongest rho, which moved
  those rows further from rho.
\end{itemize}
